\documentclass[11pt,oneside,openany,fontset=none]{ctexbook}
\usepackage{style/threeblue}

\title{最优控制与偏微分方程}
\author{TFC 课程中文精校讲义}
\date{2026}
\hypersetup{
  pdftitle={最优控制与偏微分方程：十讲中文精校课程讲义},
  pdfauthor={Wilfrid Gangbo; Andrzej Święch},
  pdfsubject={平均场博弈、Wasserstein 几何、随机最优控制、动态规划与 HJB 方程},
  pdfkeywords={最优控制, 偏微分方程, 平均场博弈, Wasserstein, HJB, 动态规划, 黏性解}
}

\begin{document}

% ==================== Cover ====================
\begin{titlepage}
\thispagestyle{empty}
\begin{tikzpicture}[remember picture,overlay]
  \fill[DeepSpace] (current page.south west) rectangle (current page.north east);
  \foreach \x in {0,...,12}{
    \draw[GridLine,opacity=.28] ($(current page.south west)+(\x*1.75cm,0)$) -- ($(current page.north west)+(\x*1.75cm,0)$);
  }
  \foreach \y in {0,...,17}{
    \draw[GridLine,opacity=.28] ($(current page.south west)+(0,\y*1.75cm)$) -- ($(current page.south east)+(0,\y*1.75cm)$);
  }
  % A family of controlled trajectories
  \draw[BlueAccent,line width=2.2pt] ($(current page.south west)+(1.8cm,6.0cm)$)
    .. controls +(3.0cm,2.1cm) and +(-4.0cm,-1.6cm) .. ($(current page.south east)+(-1.8cm,9.2cm)$);
  \draw[YellowAccent,line width=1.7pt] ($(current page.south west)+(1.8cm,6.0cm)$)
    .. controls +(3.4cm,0.2cm) and +(-3.4cm,1.7cm) .. ($(current page.south east)+(-1.8cm,9.2cm)$);
  \draw[GreenAccent,line width=1.7pt] ($(current page.south west)+(1.8cm,6.0cm)$)
    .. controls +(2.7cm,4.0cm) and +(-3.2cm,-4.0cm) .. ($(current page.south east)+(-1.8cm,9.2cm)$);
  \fill[BlueAccent] ($(current page.south west)+(1.8cm,6.0cm)$) circle (3.5pt);
  \fill[YellowAccent] ($(current page.south east)+(-1.8cm,9.2cm)$) circle (3.5pt);
  % Measure particles
  \foreach \x/\y in {2.0/20.7,2.7/21.4,3.4/20.5,4.2/21.1,5.0/20.3,5.8/21.0,6.6/20.5}
    \fill[PurpleAccent,opacity=.75] ($(current page.south west)+(\x cm,\y cm)$) circle (2.1pt);
\end{tikzpicture}

\vspace*{22mm}
\begin{flushleft}
  {\sffamily\small\bfseries\textcolor{BlueAccent}{TOHOKU FORUM FOR CREATIVITY \quad 2017}}\\[13mm]
  {\sffamily\fontsize{30}{36}\selectfont\bfseries\textcolor{SoftWhite}{最优控制}}\\[1mm]
  {\sffamily\fontsize{30}{36}\selectfont\bfseries\textcolor{SoftWhite}{与偏微分方程}}\\[7mm]
  {\sffamily\Large\textcolor{MutedText}{Optimal Control and Partial Differential Equations}}\\[11mm]
  \begin{tikzpicture}
    \fill[BlueAccent] (0,0) rectangle (3.0,.11);
    \fill[YellowAccent] (3.15,0) rectangle (5.2,.11);
    \fill[GreenAccent] (5.35,0) rectangle (7.0,.11);
  \end{tikzpicture}\\[11mm]
  {\sffamily\large\textcolor{SoftWhite}{平均场博弈 \textperiodcentered{} Wasserstein 几何 \textperiodcentered{} 动态规划 \textperiodcentered{} HJB 方程}}\\[4mm]
  {\sffamily\normalsize\textcolor{MutedText}{十讲中文精校课程讲义 · 公式推导与几何图解}}
\end{flushleft}

\vfill
\begin{flushleft}
{\sffamily\small\textcolor{MutedText}{主讲：Wilfrid Gangbo · Andrzej Święch}}\\[2mm]
{\sffamily\small\textcolor{BlueAccent}{深色屏幕阅读版 \quad 3Blue1Brown-inspired visual edition}}
\end{flushleft}
\end{titlepage}

% ==================== Front matter ====================
\frontmatter
\chapter*{使用说明}
\addcontentsline{toc}{chapter}{使用说明}

这份讲义按原课程十讲的逻辑重构，而不是逐句字幕汇编。前五讲以\keyterm{平均场博弈、Wasserstein 空间与主方程}为主线，后五讲以\keyterm{随机最优控制、动态规划原理与 HJB 方程}为主线。图形、颜色和空间关系用来表达数学结构：\blue{蓝色}表示状态或主对象，\yellow{黄色}表示价值与目标，\green{绿色}表示演化或可行方向，\red{红色}表示误区或约束。

\begin{intuition}[怎样使用这份讲义]
先看每讲开头的“本讲目标”和中间的大公式，再读几何图解；最后用“一页记忆”复述主线。若能不看正文解释每幅图中的箭头，便已掌握该讲的结构。
\end{intuition}

\begin{pitfall}[约定说明]
HJB、Fokker--Planck 与主方程在不同文献中常采用不同的时间方向、Hamiltonian 定义和扩散系数归一化。本讲义会明确当前约定；与其他资料对照时，应先检查符号，而不是只比公式外观。
\end{pitfall}

\tableofcontents

\mainmatter
\chapter*{课程地图：从一个粒子到一群理性主体}
\addcontentsline{toc}{chapter}{课程地图：从一个粒子到一群理性主体}
\markboth{课程地图}{课程地图}

\begin{learninggoals}[读完这一章，你会看到]
\begin{itemize}
  \item 两条看似独立的路线，为什么都源自 Bellman 的“向后看价值、向前走状态”；
  \item 平均场博弈、最优控制、HJB 方程、连续性方程与主方程之间的角色分工；
  \item 十讲内容的依赖关系，以及每一讲解决了哪个数学障碍。
\end{itemize}
\end{learninggoals}

\section*{一个统一问题}

考虑状态 $X_t$。我们可以选择控制 $a_t$，但未来既受动力学影响，也可能受随机噪声或其他参与者影响。优化问题的核心是
\[
\textcolor{YellowAccent}{\text{总代价最小}}
\quad\text{且}\quad
\textcolor{BlueAccent}{\text{状态服从动力学}}.
\]
如果只有一个决策者，最优未来成本由\keyterm{值函数} $V(t,x)$ 记录；如果有海量相互作用的决策者，单个主体还必须知道总体分布 $m_t$，于是价值变为 $u(t,x)$，乃至依赖整份分布的 $U(t,x,m)$。

\begin{center}
\begin{tikzpicture}[node distance=12mm and 13mm]
  \node[concept] (state) {当前状态\\$x$};
  \node[conceptyellow,right=of state] (choice) {选择控制\\$a$};
  \node[conceptgreen,right=of choice] (next) {下一时刻\\$X_{t+\Delta t}$};
  \node[conceptpurple,below=of choice] (pop) {群体分布\\$m_t$};
  \node[conceptyellow,above=of choice] (value) {未来价值\\$V$ 或 $U$};
  \draw[flow] (state)--(choice);
  \draw[flow] (choice)--(next);
  \draw[warmflow] (value)--(choice);
  \draw[softflow,<->] (pop)--(choice);
  \draw[softflow,bend left=25] (next) to (value);
\end{tikzpicture}
\captionof{figure}{最优控制的闭环：价值指导选择，选择推动状态，状态更新未来价值。平均场博弈多出群体分布这一条反馈通道。}
\end{center}

\section*{两条主线}

\begin{tcolorbox}[colback=NightPanelTwo,colframe=GridLine]
\begin{tabularx}{\textwidth}{@{}>{\sffamily\bfseries\color{BlueAccent}}p{0.18\textwidth}YY@{}}
\toprule
 & \textcolor{YellowAccent}{前五讲：群体与分布} & \textcolor{GreenAccent}{后五讲：随机控制与 DPP} \\
\midrule
对象 & 个体状态 $x$ 与总体分布 $m$ & 受控扩散状态 $X_s$ \\
几何 & Wasserstein 概率测度空间 & 欧氏空间及概率空间 \\
向后方程 & HJ / 主方程 & HJB 方程 \\
向前方程 & 连续性或 Fokker--Planck 方程 & 受控随机微分方程 \\
关键原则 & Nash 一致性、单调性、势结构 & 动态规划、验证、可测拼接 \\
弱解语言 & 测度解、黏性思想 & 黏性解与比较原理 \\
\bottomrule
\end{tabularx}
\end{tcolorbox}

\section*{十讲依赖图}

\begin{center}
\begin{tikzpicture}[node distance=8mm and 9mm,scale=.90,transform shape]
\node[concept] (l1) {01\\从 $N$ 人博弈到 MFG};
\node[concept,right=of l1] (l2) {02\\Wasserstein 几何};
\node[concept,right=of l2] (l3) {03\\Lions 导数与主方程};
\node[concept,right=of l3] (l4) {04\\特征线与流};
\node[concept,right=of l4] (l5) {05\\势平均场博弈};

\node[conceptyellow,below=14mm of l1] (l6) {06\\随机控制与 HJB};
\node[conceptyellow,right=of l6] (l7) {07\\验证与反馈};
\node[conceptyellow,right=of l7] (l8) {08\\可测结构};
\node[conceptyellow,right=of l8] (l9) {09\\DPP 严格证明};
\node[conceptyellow,right=of l9] (l10) {10\\黏性解与唯一性};

\foreach \a/\b in {l1/l2,l2/l3,l3/l4,l4/l5}{\draw[flow] (\a)--(\b);}
\foreach \a/\b in {l6/l7,l7/l8,l8/l9,l9/l10}{\draw[warmflow] (\a)--(\b);}
\draw[softflow,dashed] (l1) -- node[left,font=\scriptsize\sffamily,text=MutedText]{Bellman 思想} (l6);
\draw[softflow,dashed] (l3) -- node[right,font=\scriptsize\sffamily,text=MutedText]{无限维 HJ} (l10);
\end{tikzpicture}
\captionof{figure}{课程不是十个孤立主题，而是两条并行推进、互相照亮的路线。}
\end{center}

\section*{三组必须区分的对象}

\subsection*{状态、分布与分布上的函数}

\begin{itemize}
  \item $x\in\R^d$：某个具体个体的位置或状态；
  \item $m\in\PP_2(\R^d)$：总体状态的概率分布；
  \item $\mathcal F(m)\in\R$：给整份分布打分的泛函；
  \item $U(t,x,m)$：既看个体位置，也看总体分布的价值。
\end{itemize}

把这四类对象混在一起，是学习主方程最常见的障碍。可以把 $x$ 想成“地图上的一个点”，$m$ 是“整张热力图”，$\mathcal F$ 是“对热力图的总评分”，而 $U$ 则是“在这张热力图背景下，一个点的未来价值”。

\subsection*{开环控制与反馈控制}

开环控制 $\alpha_s$ 是预先给出的适应过程；反馈控制写成 $\alpha_s=\hat a(s,X_s)$，会根据当前状态实时调整。HJB 方程给出的通常正是最优反馈律
\[
\hat a(t,x)\in\argmin_{a\in A}
\left\{b(t,x,a)\cdot D V(t,x)
+\frac12\Tr\!\left[\sigma\sigma^\top(t,x,a)D^2V(t,x)\right]+f(t,x,a)\right\}.
\]

\subsection*{经典解与黏性解}

经典解要求偏导数真实存在；黏性解不直接对 $V$ 求二阶导，而是让光滑测试函数从上方或下方接触 $V$。这不是“更模糊”的解，而是把 HJB 的最优性结构保存到不可微情形中的正确语言。

\begin{takeaway}[课程总纲]
\begin{enumerate}
  \item \blue{状态向前演化}：SDE、连续性方程或 Fokker--Planck 方程；
  \item \yellow{价值向后传播}：HJ、HJB 或主方程；
  \item \green{最优控制把两者闭合}：价值梯度产生反馈，反馈又改变状态分布；
  \item 当光滑性消失时，\purple{动态规划 $\Rightarrow$ 黏性解 $\Rightarrow$ 比较原理}维持理论闭环。
\end{enumerate}
\end{takeaway}

\chapter{从 \texorpdfstring{$N$}{N} 人博弈到平均场博弈}
\lectureribbon{01}{On a Mean Field Game Equation I}{Wilfrid Gangbo}

\begin{learninggoals}
\begin{itemize}
  \item 从有限 $N$ 人 Nash 博弈中识别“只通过经验分布相互作用”的结构；
  \item 理解 $N\to\infty$ 后为什么出现一对前向--后向偏微分方程；
  \item 看清平均场博弈系统与主方程分别回答什么问题。
\end{itemize}
\end{learninggoals}

\section{有限玩家：每个人都在优化，但环境由所有人共同生成}

设有 $N$ 个玩家，第 $i$ 个玩家的状态为 $X_t^i\in\R^d$，控制为 $\alpha_t^i$。一个典型动力学写成
\[
\dd X_t^i=b\!\left(X_t^i,\alpha_t^i,m_t^N\right)\dd t
+\sqrt{2\nu}\,\dd W_t^i,
\qquad
m_t^N=\frac1N\sum_{j=1}^N\delta_{X_t^j}.
\]
这里的 $m_t^N$ 是\keyterm{经验测度}：它把“其他 $N-1$ 人在哪里”压缩成一份分布。玩家 $i$ 的目标为
\[
J^i(\bmalpha)=
\E\!\left[
\int_0^T L\!\left(X_t^i,\alpha_t^i,m_t^N\right)\dd t
+G\!\left(X_T^i,m_T^N\right)
\right].
\]

\begin{definition}[Nash 均衡]
控制组合 $\bmalpha^*=(\alpha^{1,*},\ldots,\alpha^{N,*})$ 是 Nash 均衡，如果任何一个玩家在其他人策略不变时都不能通过单独改策获益：
\[
J^i(\alpha^{i,*},\bmalpha^{-i,*})
\le J^i(\beta^i,\bmalpha^{-i,*})
\quad\text{对所有容许 }\beta^i.
\]
它表达的是\emterm{单边偏离无利可图}，并不等于所有人的总代价最小。
\end{definition}

\begin{center}
\begin{tikzpicture}[scale=.96]
  \draw[guide,rounded corners=8pt] (-4.6,-2.1) rectangle (4.6,2.1);
  \foreach \x/\y in {-3.5/1.1,-2.7/-.8,-1.8/.4,-.9/-1.3,.1/1.2,1.2/-.3,2.1/.8,3.1/-1.2,3.7/.25}{
    \fill[BlueAccent,opacity=.85] (\x,\y) circle (3pt);
  }
  \fill[YellowAccent] (-.9,-1.3) circle (5pt);
  \draw[warmflow] (-.9,-1.3) -- ++(1.45,.9) node[right,font=\scriptsize\sffamily,text=YellowAccent]{玩家 $i$ 的选择};
  \node[conceptpurple] at (0,-2.75) {$N$ 个点 $\longmapsto$ 经验分布 $m_t^N$};
  \draw[softflow] (0,-2.1)--(0,-2.3);
\end{tikzpicture}
\captionof{figure}{玩家只需感知群体的统计形状，而不必记住每个人的身份。黄色玩家既受分布影响，也以 $1/N$ 的权重影响分布。}
\end{center}

\section{Hamiltonian：把“选控制”压缩成一个函数}

先采用控制进入速度、运行成本为 $L(x,a)$ 的简化模型。定义 Hamiltonian
\[
H(x,p)=\sup_{a\in A}\bigl\{-a\cdot p-L(x,a)\bigr\}.
\]
符号选择取决于将问题写成最小化还是最大化。本讲义采用上式时，最优速度满足
\[
a^*(x,p)\in\argmax_{a\in A}\{-a\cdot p-L(x,a)\},
\qquad a^*=-D_pH(x,p)
\]
（后一式需 $H$ 光滑且采用对应约定）。因此，价值梯度 $p=\nabla u$ 告诉玩家“往哪个方向移动最能降低未来总代价”。

\begin{examplebox}[二次动能]
若 $L(a)=\tfrac12\abs a^2$，则
\[
H(p)=\sup_a\{-a\cdot p-\tfrac12\abs a^2\}
=\frac12\abs p^2,
\qquad a^*(p)=-p.
\]
价值上升最快的方向是 $\nabla u$，所以最优运动朝 $-\nabla u$。
\end{examplebox}

\section{极限 \texorpdfstring{$N\to\infty$}{N→∞}：一个代表性玩家面对确定分布流}

当玩家可交换，且每个人对经验分布的影响只有 $1/N$ 时，极限中单个玩家近似把总体分布流 $m=(m_t)_{0\le t\le T}$ 当成给定环境。固定 $m$，代表性玩家求解 HJB 方程；所有玩家采用该最优反馈后，其状态分布又必须恰好等于 $m$。这就是\keyterm{平均场一致性条件}。

在常见扩散约定下，平均场博弈系统为
\focusequation{平均场博弈的前向--后向系统}{%
\[
\left\{
\begin{aligned}
-\partial_t u-\nu\Delta u+H(x,\nabla u)&=F(x,m_t),
&&u(T,x)=G(x,m_T),\\
\partial_t m-\nu\Delta m-\nabla\!\cdot\!\bigl(mD_pH(x,\nabla u)\bigr)&=0,
&&m(0)=m_0.
\end{aligned}
\right.
\]}

第一行\yellow{从终端向过去}传播价值，第二行\green{从初始时刻向未来}传播人口。两行通过 $F,G$ 与 $D_pH$ 双向耦合。

\begin{center}
\begin{tikzpicture}[node distance=16mm and 20mm]
  \node[conceptyellow] (u) {价值 $u(t,x)$\\向后求解};
  \node[conceptgreen,right=of u] (m) {分布 $m_t$\\向前演化};
  \draw[warmflow,bend left=18] (u) to node[above,font=\scriptsize\sffamily,text=MutedText]{最优反馈 $-D_pH(x,\nabla u)$} (m);
  \draw[flow,bend left=18] (m) to node[below,font=\scriptsize\sffamily,text=MutedText]{拥挤成本 $F(x,m_t)$} (u);
  \node[below=13mm of $(u)!0.5!(m)$,font=\sffamily\small,text=BlueAccent] {不动点：猜测的分布 = 最优行为产生的分布};
\end{tikzpicture}
\captionof{figure}{MFG 的核心不是单个 PDE，而是价值与分布之间的闭环不动点。}
\end{center}

\section{为什么还需要主方程}

MFG 系统针对一个给定初始分布 $m_0$，求出整条均衡轨道 $(u_t,m_t)$。但如果我们想同时描述“从任意时刻、任意个体状态、任意总体分布出发”的均衡价值，就要引入
\[
U(t,x,m): [0,T]\times\R^d\times\PP_2(\R^d)\longrightarrow\R.
\]
沿均衡分布轨道应有 $u(t,x)=U(t,x,m_t)$。对这个复合函数使用分布方向的链式法则，就会得到\keyterm{主方程}。因此：

\begin{itemize}
  \item MFG 系统描述一条均衡轨道；
  \item 主方程描述所有可能分布状态上的均衡反馈场；
  \item 有了足够光滑的 $U$，可用 $U(t,x,m^N)$ 近似有限 $N$ 人博弈的值函数。
\end{itemize}

\begin{pitfall}[社会规划问题不等于 Nash 博弈]
把所有玩家成本相加后统一最小化，得到的是社会规划者的最优控制；Nash 均衡只要求每个人对单边偏离最优。只有在特殊的势结构下，两者才可由同一个整体泛函联系起来，第五讲将回到这一点。
\end{pitfall}

\section{一个静态拥挤例子}

设终端成本 $G(x,m)=g(x)+\kappa\rho_m(x)$，其中 $\rho_m$ 是分布密度，$\kappa>0$。若所有人涌向 $g$ 最小的地点，该处密度升高，拥挤项又会抬高成本。均衡不是“所有人去同一点”，而是外部吸引 $g$ 与内生排斥 $\kappa\rho_m$ 的平衡。

\begin{center}
\begin{tikzpicture}[x=1cm,y=1cm]
  \draw[axis] (-4.2,0)--(4.2,0) node[right,font=\scriptsize\sffamily]{位置 $x$};
  \draw[axis] (0,-.2)--(0,2.9) node[above,font=\scriptsize\sffamily]{成本};
  \draw[BlueAccent,very thick,domain=-3.6:3.6,samples=80] plot (\x,{0.18*(\x)^2+.25});
  \draw[RedAccent,very thick,domain=-3.6:3.6,samples=80] plot (\x,{2.25*exp(-0.75*(\x)^2)});
  \draw[YellowAccent,line width=2.2pt,domain=-3.6:3.6,samples=80] plot (\x,{0.18*(\x)^2+.25+2.25*exp(-0.75*(\x)^2)});
  \node[text=BlueAccent,font=\scriptsize\sffamily] at (3.15,2.05) {外部成本 $g$};
  \node[text=RedAccent,font=\scriptsize\sffamily] at (-2.65,1.85) {拥挤成本};
  \node[text=YellowAccent,font=\scriptsize\sffamily] at (2.6,.85) {总成本};
\end{tikzpicture}
\captionof{figure}{个体最优与群体分布相互塑造：分布不是参数，而是均衡的一部分。}
\end{center}

\begin{takeaway}
\begin{enumerate}
  \item 经验测度 $m_t^N$ 把对称的 $N$ 人相互作用压缩成分布；
  \item $N\to\infty$ 后，代表性玩家对给定 $m$ 做最优控制，再以一致性条件闭环；
  \item MFG 是“向后价值 + 向前分布”的耦合系统；
  \item 主方程把价值提升为 $U(t,x,m)$，统一编码所有初始分布。
\end{enumerate}
\end{takeaway}

\begin{exercisebox}
\begin{enumerate}
  \item 对 $L(a)=\tfrac{1}{q}\abs a^q$（$q>1$），计算对应 Hamiltonian 的幂指数，并解释最优速度方向。
  \item 若 $F(x,m)$ 与 $m$ 无关，MFG 系统还保留哪些耦合？它是否退化为普通最优控制？
  \item 用一句话区分：经验测度、总体极限分布、主方程中的分布变量。
\end{enumerate}
\end{exercisebox}

\chapter{Wasserstein 空间的几何}
\lectureribbon{02}{On a Mean Field Game Equation II}{Wilfrid Gangbo}
\begin{learninggoals}
\begin{itemize}
  \item 把“两个分布有多远”写成最优搬运问题；
  \item 理解概率测度空间为何像一个无限维流形：有曲线、速度和切空间；
  \item 掌握沿连续性方程的链式法则，为主方程中的分布导数做准备。
\end{itemize}
\end{learninggoals}

\section{概率分布为什么需要自己的距离}

平均场博弈中的状态不只是 $x\in\R^d$，还包括一份概率分布 $m$。弱收敛能够描述“积分是否收敛”，却没有直接告诉我们移动质量需要多少代价。对于有限二阶矩的概率测度
\[
\PP_2(\R^d)=\left\{\mu:\ \mu(\R^d)=1,
\ \int_{\R^d}\abs x^2\mu(\dd x)<\infty\right\},
\]
二阶 Wasserstein 距离把几何位置纳入比较。

\begin{definition}[耦合与 Wasserstein 距离]
若 $\gamma$ 在 $\R^d\times\R^d$ 上的两个边缘分布分别是 $\mu,\nu$，则称 $\gamma\in\Gamma(\mu,\nu)$ 是二者的耦合。定义
\focusequation{$W_2$：最省力的搬运方案}{%
\[
W_2^2(\mu,\nu)=
\inf_{\gamma\in\Gamma(\mu,\nu)}
\int_{\R^d\times\R^d}\abs{x-y}^2\,\gamma(\dd x,\dd y).
\]}
耦合 $\gamma$ 说明从 $x$ 出发的质量有多少被送到 $y$；积分是总搬运能量。
\end{definition}

\begin{center}
\begin{tikzpicture}[scale=.95]
  \node[font=\sffamily\small,text=BlueAccent] at (-3.3,2.2) {源分布 $\mu$};
  \node[font=\sffamily\small,text=YellowAccent] at (3.3,2.2) {目标分布 $\nu$};
  \foreach \x/\y in {-4/1.3,-3.5/.4,-3.1/1.6,-2.7/-.5,-2.2/.7}{\fill[BlueAccent] (\x,\y) circle (3pt);}
  \foreach \x/\y in {2.1/1.1,2.8/-.3,3.25/1.6,3.75/.25,4.2/1.15}{\fill[YellowAccent] (\x,\y) circle (3pt);}
  \draw[flow,opacity=.85] (-4,1.3) to[bend left=14] (2.1,1.1);
  \draw[flow,opacity=.65] (-3.5,.4) to[bend right=8] (2.8,-.3);
  \draw[flow,opacity=.75] (-3.1,1.6) to[bend left=6] (3.25,1.6);
  \draw[flow,opacity=.55] (-2.7,-.5) to[bend right=15] (3.75,.25);
  \draw[flow,opacity=.65] (-2.2,.7) to[bend left=10] (4.2,1.15);
  \node[conceptpurple] at (0,-1.65) {在所有连线方案中，选择平方距离总和最小者};
\end{tikzpicture}
\captionof{figure}{Wasserstein 距离不是比较柱状图的高度，而是计算把一团质量搬成另一团质量所需的最小功。}
\end{center}

若最优耦合由映射 $T$ 诱导，即 $T_\#\mu=\nu$，则
\[
W_2^2(\mu,\nu)=\int\abs{x-T(x)}^2\mu(\dd x).
\]
这里 $T_\#\mu$ 称为\keyterm{推前测度}，含义是：若 $X\sim\mu$，则 $T(X)\sim T_\#\mu$。

\begin{examplebox}[两个 Dirac 质量]
对 $\mu=\delta_a,\nu=\delta_b$，只有一种耦合，故 $W_2(\delta_a,\delta_b)=\abs{a-b}$。这说明 Wasserstein 空间把欧氏空间作为 Dirac 测度子空间嵌入其中。
\end{examplebox}

\section{从随机变量看：Wasserstein 空间是 Hilbert 空间的商}

取一个足够丰富的无原子概率空间 $(\Omega,\mathcal A,\Pp)$，令
\[
\mathcal H=L^2(\Omega;\R^d),\qquad X\longmapsto\Law(X).
\]
许多不同的随机变量具有同一分布；因此 law map 把 $\mathcal H$ 中所有“仅差重排”的点识别起来。关键恒等式是
\[
W_2^2(\mu,\nu)
=\inf\left\{\E\abs{X-Y}^2:\ \Law(X)=\mu,\ \Law(Y)=\nu\right\}.
\]

\begin{intuition}[商空间图像]
Hilbert 空间里有大量带标签的粒子配置。抹去标签，只保留直方图或概率云，就得到 Wasserstein 空间。$W_2$ 是在所有可能标签匹配中选出的最短 $L^2$ 距离。
\end{intuition}

\begin{center}
\begin{tikzpicture}[node distance=15mm]
  \node[concept] (x1) {$X_1(\omega)$};
  \node[concept,below=7mm of x1] (x2) {$X_2(\omega)$};
  \node[concept,below=7mm of x2] (x3) {$X_3(\omega)$};
  \node[conceptpurple,right=28mm of x2,minimum height=22mm] (law) {同一分布\\$\mu=\Law(X_k)$};
  \draw[flow] (x1)--(law);
  \draw[flow] (x2)--(law);
  \draw[flow] (x3)--(law);
  \node[font=\scriptsize\sffamily,text=MutedText,below=4mm of law] {抹去样本标签，只保留质量位置};
\end{tikzpicture}
\captionof{figure}{law map 的纤维：不同随机变量可以代表 Wasserstein 空间中的同一点。}
\end{center}

\section{测度曲线、速度场与连续性方程}

在欧氏空间中，曲线 $x_t$ 的速度是 $\dot x_t$。在概率测度空间中，曲线 $\mu_t$ 的速度由向量场 $v_t(x)$ 表示，并满足
\focusequation{质量守恒}{%
\[
\partial_t\mu_t+\nabla\cdot(v_t\mu_t)=0.
\]}
这就是连续性方程。它的弱形式为：对任意光滑紧支撑测试函数 $\varphi$，
\[
\frac{\dd}{\dd t}\int\varphi(x)\mu_t(\dd x)
=\int\nabla\varphi(x)\cdot v_t(x)\mu_t(\dd x).
\]
速度的最小动能 $\norm{v_t}_{L^2(\mu_t)}$ 等于曲线的度量速度 $\abs{\dot\mu_t}$。因此“测度在动”可以严格地翻译成“质量被速度场搬运”。

\begin{definition}[Wasserstein 切空间]
在 $\mu$ 处的切空间通常取为梯度场在 $L^2(\mu)$ 中的闭包：
\[
T_\mu\PP_2(\R^d)
=\overline{\{\nabla\phi:\ \phi\in C_c^\infty(\R^d)\}}^{L^2(\mu)}.
\]
只保留梯度部分，是因为同一个测度曲线可由多个速度场表示；最小范数代表是梯度场。
\end{definition}

\section{Wasserstein 测地线}

若 $T$ 是从 $\mu_0$ 到 $\mu_1$ 的最优搬运映射，则
\[
\mu_s=\bigl((1-s)\operatorname{Id}+sT\bigr)_\#\mu_0,
\qquad 0\le s\le1,
\]
是一条常速测地线。每个质量粒子沿直线前进，但整份分布走的是 Wasserstein 空间中的“直线”。

\begin{center}
\begin{tikzpicture}[x=1.1cm,y=1cm]
  \draw[softflow] (-4,0)--(4.2,0) node[right,font=\scriptsize\sffamily] {$s$};
  \foreach \s/\c in {0/BlueAccent,1.5/PurpleAccent,3/GreenAccent,4.5/YellowAccent}{
    \draw[\c,very thick,domain=-1.2:1.2,samples=50] plot ({\s-2.25+.18*\x},{1.55*exp(-2.2*(\x)^2)-1.15});
  }
  \node[font=\scriptsize\sffamily,text=BlueAccent] at (-2.25,.75) {$\mu_0$};
  \node[font=\scriptsize\sffamily,text=YellowAccent] at (2.25,.75) {$\mu_1$};
  \draw[flow] (-1.9,1.2)--(1.9,1.2);
\end{tikzpicture}
\captionof{figure}{测地线把最优匹配的每一小团质量按恒定速度插值。}
\end{center}

\section{分布泛函的梯度与链式法则}

设 $\mathcal F:\PP_2(\R^d)\to\R$ 足够光滑。它的 Wasserstein 梯度记为 $D_\mu\mathcal F(\mu)(x)\in T_\mu\PP_2$。若 $(\mu_t,v_t)$ 满足连续性方程，则
\focusequation{分布空间的链式法则}{%
\[
\frac{\dd}{\dd t}\mathcal F(\mu_t)
=\int_{\R^d}
D_\mu\mathcal F(\mu_t)(x)\cdot v_t(x)\,\mu_t(\dd x).
\]}
它与欧氏链式法则 $\frac{\dd}{\dd t}f(x_t)=\nabla f(x_t)\cdot\dot x_t$ 完全同构，只是内积变成了 $L^2(\mu_t)$ 内积。

\begin{examplebox}[二阶矩泛函]
令 $\mathcal F(\mu)=\tfrac12\int\abs x^2\mu(\dd x)$，则 $D_\mu\mathcal F(\mu)(x)=x$。于是
\[
\frac{\dd}{\dd t}\mathcal F(\mu_t)=\int x\cdot v_t(x)\mu_t(\dd x).
\]
若 $v_t(x)=-x$，二阶矩以指数速度下降，质量向原点收缩。
\end{examplebox}

\begin{pitfall}[“对密度求导”不是完整定义]
测度可能没有密度，或者含有 Dirac 质量。因此 Wasserstein 微分必须以推前、耦合或 Hilbert 提升为基础，而不能只对 $\rho(x)$ 做形式偏导。下一讲的 Lions 导数正是把这一点变成可计算工具。
\end{pitfall}

\begin{takeaway}
\begin{enumerate}
  \item $W_2$ 衡量最小二次搬运成本；
  \item 随机变量的 $L^2$ 空间经“同分布”识别后，产生 Wasserstein 几何；
  \item 测度曲线的速度由连续性方程中的 $v_t$ 表示；
  \item 分布泛函的链式法则是普通梯度内积的无限维版本。
\end{enumerate}
\end{takeaway}

\begin{exercisebox}
\begin{enumerate}
  \item 证明 $W_2(\delta_a,\delta_b)=\abs{a-b}$，并写出对应测地线。
  \item 若 $\mu_t=(e^{-t}\operatorname{Id})_\#\mu_0$，求一个满足连续性方程的速度场。
  \item 对 $\mathcal F(\mu)=\int V(x)\mu(\dd x)$，猜测 $D_\mu\mathcal F$，再用链式法则验证。
\end{enumerate}
\end{exercisebox}

\chapter{Lions 导数与主方程}
\lectureribbon{03}{On a Mean Field Game Equation III}{Wilfrid Gangbo}
\begin{learninggoals}
\begin{itemize}
  \item 用 Hilbert 提升把“对概率测度求导”化为 Fréchet 导数；
  \item 区分线性泛函导数 $\delta/\delta m$ 与 Lions 导数 $D_m$；
  \item 从均衡流上的链式法则看懂主方程的每一项；
  \item 理解主方程如何统一近似有限玩家值函数。
\end{itemize}
\end{learninggoals}

\section{Hilbert 提升：先给随机变量求导，再忘掉标签}

设 $\mathcal F:\PP_2(\R^d)\to\R$。在 $\mathcal H=L^2(\Omega;\R^d)$ 上定义提升
\[
\widetilde{\mathcal F}(X)=\mathcal F(\Law(X)).
\]
提升函数对所有保持分布的重排不变。若 $\widetilde{\mathcal F}$ 在 $X$ 处 Fréchet 可微，则存在 $D\widetilde{\mathcal F}(X)\in\mathcal H$，使
\[
\widetilde{\mathcal F}(X+Y)-\widetilde{\mathcal F}(X)
=\E\!\left[D\widetilde{\mathcal F}(X)\cdot Y\right]+o(\norm Y_{L^2}).
\]
重排不变性迫使该导数只通过 $X$ 的取值发挥作用：存在向量场 $D_m\mathcal F(\mu)(\cdot)$，使得
\focusequation{Lions 导数的表示}{%
\[
D\widetilde{\mathcal F}(X)
=D_m\mathcal F(\Law(X))(X).
\]}

\begin{center}
\begin{tikzpicture}[node distance=14mm and 18mm]
  \node[concept] (measure) {$\mathcal F(\mu)$\\测度空间};
  \node[conceptpurple,right=of measure] (lift) {$\widetilde{\mathcal F}(X)$\\Hilbert 空间};
  \node[conceptyellow,right=of lift] (frechet) {$D\widetilde{\mathcal F}(X)$\\Fréchet 导数};
  \node[conceptgreen,below=of lift] (lions) {$D_m\mathcal F(\mu)(x)$\\Lions 导数};
  \draw[flow] (measure)--node[above,font=\scriptsize\sffamily,text=MutedText]{提升} (lift);
  \draw[warmflow] (lift)--node[above,font=\scriptsize\sffamily,text=MutedText]{求导} (frechet);
  \draw[softflow] (frechet) to[bend left=10] node[right,font=\scriptsize\sffamily,text=MutedText]{按 $X$ 表示} (lions);
  \draw[softflow] (lions) to[bend left=10] node[left,font=\scriptsize\sffamily,text=MutedText]{下降到分布} (measure);
\end{tikzpicture}
\captionof{figure}{Lions 微分的策略：在平坦的 Hilbert 空间求导，再利用对称性降回概率测度空间。}
\end{center}

\section{两种常见“分布导数”}

若存在标量函数 $\frac{\delta\mathcal F}{\delta m}(m,x)$，使小扰动 $m+\varepsilon(\nu-m)$ 满足
\[
\left.\frac{\dd}{\dd\varepsilon}\mathcal F\bigl(m+\varepsilon(\nu-m)\bigr)\right|_{\varepsilon=0}
=\int\frac{\delta\mathcal F}{\delta m}(m,x)(\nu-m)(\dd x),
\]
则它称为线性泛函导数，只确定到一个与 $x$ 无关的常数。Wasserstein/Lions 向量导数通常是
\[
D_m\mathcal F(m,x)=D_x\frac{\delta\mathcal F}{\delta m}(m,x).
\]

\begin{examplebox}[势能与相互作用能]
若
\[
\mathcal F(m)=\int V(x)m(\dd x)
+\frac12\iint W(x-y)m(\dd x)m(\dd y),
\]
则
\[
\frac{\delta\mathcal F}{\delta m}(m,x)
=V(x)+(W*m)(x),\qquad
D_m\mathcal F(m,x)=\nabla V(x)+(\nabla W*m)(x).
\]
第一项是外场，第二项是群体平均相互作用力。
\end{examplebox}

\section{主方程的来源：沿所有均衡流同时使用链式法则}

令 $U(t,x,m)$ 表示在时刻 $t$、个体状态 $x$、群体分布 $m$ 下的均衡价值。对某条均衡流 $m_t$，令
\[
u(t,x)=U(t,x,m_t).
\]
那么
\[
\frac{\dd}{\dd t}U(t,x,m_t)
=\partial_tU(t,x,m_t)
+\int D_mU(t,x,m_t,y)\cdot v_t(y)m_t(\dd y).
\]
当群体最优速度为 $v_t(y)=-D_pH(y,D_xU(t,y,m_t))$ 时，把这条链式法则代入个体 HJ 方程，就得到第一阶主方程的典型结构：

\begin{derivation}[第一阶主方程：结构式]
采用第一讲的 Hamiltonian 约定，一种常见写法是
\[
\begin{aligned}
-\partial_tU(t,x,m)
&+H\bigl(x,D_xU(t,x,m)\bigr)\\
&+\int_{\R^d}D_mU(t,x,m,y)\cdot
D_pH\bigl(y,D_xU(t,y,m)\bigr)m(\dd y)
=F(x,m),
\end{aligned}
\]
终端条件为 $U(T,x,m)=G(x,m)$。加入个体噪声或共同噪声后，还会出现关于 $x$ 与 $m$ 的二阶项；不同 Hamiltonian 定义会改变若干符号。
\end{derivation}

\begin{center}
\begin{tikzpicture}[node distance=12mm and 15mm]
  \node[conceptyellow] (time) {$-\partial_tU$\\时间回溯};
  \node[concept,right=of time] (self) {$H(x,D_xU)$\\个体最优};
  \node[conceptpurple,right=of self] (mean) {$\int D_mU\cdot v\,\dd m$\\分布被搬运};
  \node[conceptgreen,right=of mean] (cost) {$F(x,m)$\\群体成本};
  \draw[warmflow] (time)--(self);
  \draw[flow] (self)--(mean);
  \draw[softflow] (mean)--(cost);
\end{tikzpicture}
\captionof{figure}{主方程的四块拼图：时间、个体状态、总体分布与耦合成本。}
\end{center}

\section{从主方程回到 MFG 系统}

若 $U$ 足够光滑，给定 $m_0$，先解前向方程
\[
\partial_t m_t-\nabla\cdot\left(m_tD_pH\bigl(x,D_xU(t,x,m_t)\bigr)\right)=0,
\qquad m(0)=m_0,
\]
再令 $u(t,x)=U(t,x,m_t)$。链式法则保证 $(u,m)$ 满足 MFG 系统。换句话说，主方程是一台“均衡解生成器”：输入任意 $m_0$，输出对应的均衡轨道。

\section{有限玩家近似：一个函数编码所有玩家}

对称 $N$ 人博弈中，第 $i$ 个玩家的值函数依赖 $(x_1,\ldots,x_N)$。主方程建议近似
\[
v^{N,i}(t,x_1,\ldots,x_N)
\approx U\!\left(t,x_i,m_{\symbf{x}}^{N,i}\right),
\qquad
m_{\symbf{x}}^{N,i}=\frac1{N-1}\sum_{j\ne i}\delta_{x_j}.
\]
对 $x_j$ 求导时，分布变量的变化带来 $1/(N-1)$ 因子：
\[
D_{x_j}v^{N,i}
\approx\frac1{N-1}D_mU\!\left(t,x_i,m_{\symbf{x}}^{N,i},x_j\right),
\qquad j\ne i.
\]
这正是主方程能近似 Nash 系统的尺度机制：每个其他玩家的影响很小，但 $N-1$ 个小影响加总为 $O(1)$。

\begin{intuition}[为什么常先得到短时经典解]
主方程是高维、非线性、带分布变量的特征方程。短时间内，特征流通常尚未折叠，导数估计可闭合；时间太长时，不同特征可能交叉，经典光滑性会丢失。因此“短时存在”不是技术偶然，而与 HJ 方程形成激波的几何现象一致。
\end{intuition}

\section{解流与半群记号}

令 $\Sigma_s^t$ 表示从时刻 $s$ 的分布出发，沿均衡反馈演化到时刻 $t$：
\[
m_t=\Sigma_s^t(m_s),\qquad
\Sigma_r^t=\Sigma_s^t\circ\Sigma_r^s
\quad(r\le s\le t).
\]
半群性质就是“先走到 $s$，再从 $s$ 走到 $t$，等于一次走完”。它把均衡轨道组织成流，也是下一讲特征线方法的核心。

\begin{pitfall}[主方程不是把 $m$ 当有限维参数]
$m$ 的扰动是一整个位移场，导数 $D_mU(t,x,m,y)$ 还多出一个空间变量 $y$。漏掉 $y$，就看不见群体在不同位置移动对个体价值的不同影响。
\end{pitfall}

\begin{takeaway}
\begin{enumerate}
  \item Lions 导数来自 law-invariant Hilbert 提升的 Fréchet 导数；
  \item $D_m$ 是向量场，而 $\delta/\delta m$ 是只确定到常数的标量；
  \item 主方程是沿均衡分布流对 $U(t,x,m_t)$ 使用链式法则的结果；
  \item $U(t,x,m)$ 同时生成 MFG 轨道，并以 $1/N$ 尺度近似有限玩家 Nash 系统。
\end{enumerate}
\end{takeaway}

\begin{exercisebox}
\begin{enumerate}
  \item 对 $\mathcal F(m)=\abs{\int x\,m(\dd x)}^2$，计算 $\delta\mathcal F/\delta m$ 与 $D_m\mathcal F$。
  \item 解释为什么 $D_{x_j}U(t,x_i,m_{\symbf{x}}^{N,i})$ 自然带 $1/(N-1)$。
  \item 沿给定连续性方程重新推导主方程中的积分项，并核对符号。
\end{enumerate}
\end{exercisebox}

\chapter{特征线、连续性方程与流}
\lectureribbon{04}{On a Mean Field Game Equation IV}{Wilfrid Gangbo}
\begin{learninggoals}
\begin{itemize}
  \item 在“跟着粒子走”的 Lagrangian 描述与“站在空间看密度”的 Eulerian 描述之间切换；
  \item 从特征线推出连续性方程，并理解流的半群与可逆性；
  \item 看懂经典主方程解如何沿特征流构造，以及光滑性为何会失效。
\end{itemize}
\end{learninggoals}

\section{两种观察同一群体的方式}

设初始标签为 $a$，粒子轨迹 $X_t(a)$ 满足
\[
\dot X_t(a)=v_t(X_t(a)),\qquad X_0(a)=a.
\]
若初始标签按 $m_0$ 分布，则时刻 $t$ 的分布是
\[
m_t=(X_t)_\#m_0.
\]
左式跟踪每个标签，是\keyterm{Lagrangian 描述}；右式只记录每个位置有多少质量，是\keyterm{Eulerian 描述}。

\begin{center}
\begin{tikzpicture}[scale=.95]
  \node[concept] (lag) at (-3.5,2.2) {Lagrangian\\粒子轨迹 $X_t(a)$};
  \node[conceptgreen] (eul) at (3.5,2.2) {Eulerian\\密度/测度 $m_t$};
  \draw[flow,<->] (lag)--node[above,font=\scriptsize\sffamily,text=MutedText]{$m_t=(X_t)_\#m_0$} (eul);
  \foreach \y/\c in {-1.2/BlueAccent,-.45/PurpleAccent,.3/GreenAccent,1.05/YellowAccent}{
    \draw[\c,thick] (-4.4,-1.2+\y*.15) .. controls (-1.5,\y) and (1.4,-\y*.25) .. (4.3,\y*.55);
  }
  \draw[axis] (-4.6,-1.7)--(4.6,-1.7) node[right,font=\scriptsize\sffamily] {$t$};
  \node[font=\scriptsize\sffamily,text=MutedText] at (0,-2.15) {每条曲线是一粒质量；任意竖切片给出该时刻的总体分布};
\end{tikzpicture}
\captionof{figure}{轨迹与密度不是两套动力学，而是同一质量运动的微观和宏观视角。}
\end{center}

\section{从推前到连续性方程}

对测试函数 $\varphi$，由推前定义
\[
\int\varphi(x)m_t(\dd x)=\int\varphi(X_t(a))m_0(\dd a).
\]
对时间求导：
\[
\frac{\dd}{\dd t}\int\varphi\dd m_t
=\int\nabla\varphi(X_t(a))\cdot\dot X_t(a)m_0(\dd a)
=\int\nabla\varphi(x)\cdot v_t(x)m_t(\dd x).
\]
分部积分便得到
\focusequation{Eulerian 质量守恒}{%
\[
\partial_tm_t+\nabla\cdot(m_tv_t)=0.
\]}

\begin{intuition}[散度项的含义]
对一个小区域，$\partial_tm$ 记录内部质量变化，$\nabla\cdot(mv)$ 记录净流出。二者之和为零，表示粒子没有凭空产生或消失。
\end{intuition}

\section{HJ 方程的特征线：价值梯度变成共轭动量}

考虑一阶 Hamilton--Jacobi 方程
\[
-\partial_tu(t,x)+H(x,D_xu(t,x))=F(x,m_t).
\]
令共轭动量 $P_t=D_xu(t,X_t)$。在光滑情形下，特征系统具有 Hamilton 结构：
\[
\dot X_t=-D_pH(X_t,P_t),
\qquad
\dot P_t=D_xH(X_t,P_t)-D_xF(X_t,m_t),
\]
符号随 HJ 约定相应改变。$X$ 告诉质量向哪里走，$P$ 告诉价值曲面在该处朝哪里倾斜。

\begin{examplebox}[二次 Hamiltonian 与 Newton 定律]
若 $H(x,p)=\tfrac12\abs p^2-V(x)$，忽略耦合项并采用 $\dot X=-P$ 的约定，则
\[
\ddot X=-\dot P=\nabla V(X).
\]
调整 Hamiltonian 的正负约定后便得到熟悉的 $\ddot X=-\nabla V(X)$。要点不是某个固定负号，而是 HJ 特征线等价于位置--动量的 Hamilton 动力学。
\end{examplebox}

\section{流、逆流与半群}

记 $\Phi_s^t(a)$ 为从时刻 $s$、位置 $a$ 出发的特征轨迹在时刻 $t$ 的位置。若速度场足够光滑，唯一性给出
\[
\Phi_s^t\circ\Phi_r^s=\Phi_r^t,
\qquad
(\Phi_s^t)^{-1}=\Phi_t^s.
\]
分布流相应满足
\[
m_t=(\Phi_s^t)_\#m_s,
\qquad
\Sigma_s^t(m_s):=(\Phi_s^t)_\#m_s,
\qquad
\Sigma_s^t\circ\Sigma_r^s=\Sigma_r^t.
\]

\begin{center}
\begin{tikzpicture}[node distance=16mm]
  \node[concept] (r) {$m_r$};
  \node[conceptpurple,right=of r] (s) {$m_s=\Sigma_r^s(m_r)$};
  \node[conceptgreen,right=of s] (t) {$m_t=\Sigma_s^t(m_s)$};
  \draw[flow] (r)--(s);
  \draw[flow] (s)--(t);
  \draw[warmflow,bend right=33] (r) to node[below,font=\scriptsize\sffamily,text=YellowAccent]{$\Sigma_r^t$} (t);
\end{tikzpicture}
\captionof{figure}{半群性质是动力系统的“可拼接性”：时间分段不改变最终结果。}
\end{center}

\section{价值函数作为最小作用量}

对给定分布流 $m_t$，从 $(s,x)$ 出发的确定性价值可写为
\[
u(s,x)=\inf_{\substack{\gamma(s)=x\\\dot\gamma\ \text{容许}}}
\left\{
\int_s^T\bigl[L(\gamma_t,\dot\gamma_t)+F(\gamma_t,m_t)\bigr]\dd t
+G(\gamma_T,m_T)
\right\}.
\]
最优曲线满足 Euler--Lagrange/Hamilton 方程；另一方面，Bellman 原理给出 HJ 方程。\emterm{变分、特征线和 PDE} 是同一优化结构的三种语言。

\begin{derivation}[沿特征线恢复价值]
若 $u$ 光滑且 $X_t$ 采用最优速度，则
\[
\frac{\dd}{\dd t}u(t,X_t)
=\partial_tu+D_xu\cdot\dot X_t.
\]
利用 HJ 方程与 Legendre 对偶，右端化为负的最优运行成本。积分到 $T$，便恢复上面的最小作用量公式。这也给出经典解的验证方法。
\end{derivation}

\section{从流构造主方程解}

给定任意三元组 $(t,x,m)$，以 $m$ 为初始分布启动均衡特征流 $m_s=\Sigma_t^s(m)$，再计算代表性玩家从 $x$ 出发的价值，便定义 $U(t,x,m)$。如果流对初值 $m$ 光滑依赖，$U$ 的 Lions 导数存在，第三讲的链式法则就把这个构造验证为主方程经典解。

\begin{pitfall}[特征线交叉就是光滑性危机]
若两个不同初始点被流送到同一点，$\Phi_s^t$ 不再可逆；价值函数的梯度可能形成折角。此后经典特征线不能覆盖全局，需要弱解或黏性解。第十讲会用测试函数取代不存在的导数。
\end{pitfall}

\begin{examplebox}[线性流]
若 $v_t(x)=Ax$，则 $\Phi_0^t(x)=e^{tA}x$，$m_t=(e^{tA})_\#m_0$。只要矩阵指数可逆，流满足半群；若 $A=-I$，质量指数收缩且二阶矩按 $e^{-2t}$ 衰减。
\end{examplebox}

\begin{takeaway}
\begin{enumerate}
  \item $X_t$ 跟踪粒子，$m_t=(X_t)_\#m_0$ 跟踪分布；
  \item 推前公式对测试函数求导，得到连续性方程；
  \item HJ 的特征线是位置--动量 Hamilton 系统；
  \item 流的半群支持动态拼接，流的折叠预示经典解失效。
\end{enumerate}
\end{takeaway}

\begin{exercisebox}
\begin{enumerate}
  \item 从 $m_t=(\Phi_0^t)_\#m_0$ 严格写出连续性方程的弱形式。
  \item 对常速度场 $v(x)=c$ 求 $\Phi_s^t$ 与 $m_t$，并验证半群性质。
  \item 说明“最优轨迹满足特征方程”与“值函数满足 HJ 方程”之间的 Legendre 对偶步骤。
\end{enumerate}
\end{exercisebox}

\chapter{势平均场博弈与无限维 Hamilton--Jacobi 方程}
\lectureribbon{05}{On a Mean Field Game Equation V}{Wilfrid Gangbo}
\begin{learninggoals}
\begin{itemize}
  \item 识别耦合项何时来自分布泛函的一阶变分，即“势结构”；
  \item 把势平均场博弈写成 Wasserstein 空间上的整体变分问题；
  \item 理解经验测度近似中的 $1/N$ 与 $1/N^2$ 导数尺度；
  \item 看见有限维 Hamilton--Jacobi 方程与 Hilbert/Wasserstein HJ 方程的对应。
\end{itemize}
\end{learninggoals}

\section{势结构：许多个体的一阶条件来自同一个能量}

若存在分布泛函 $\mathcal F,\mathcal G$，使耦合成本满足
\[
F(x,m)=\frac{\delta\mathcal F}{\delta m}(m,x),
\qquad
G(x,m)=\frac{\delta\mathcal G}{\delta m}(m,x),
\]
则称该 MFG 具有\keyterm{势结构}。这与有限维势博弈相同：每个玩家看到的边际成本，都是同一个全局势能沿该玩家方向的导数。

\begin{center}
\begin{tikzpicture}[node distance=13mm and 20mm]
  \node[conceptpurple] (energy) {整体势能\\$\mathcal F(m)$};
  \node[conceptyellow,below left=of energy] (p1) {玩家 1\\边际成本};
  \node[conceptyellow,below=of energy] (p2) {玩家 $i$\\边际成本};
  \node[conceptyellow,below right=of energy] (pn) {玩家 $N$\\边际成本};
  \draw[flow] (energy)--node[left,font=\scriptsize\sffamily,text=MutedText]{变分} (p1);
  \draw[flow] (energy)--(p2);
  \draw[flow] (energy)--node[right,font=\scriptsize\sffamily,text=MutedText]{变分} (pn);
\end{tikzpicture}
\captionof{figure}{势结构把许多彼此耦合的一阶最优性条件整合为一个全局能量的变分。}
\end{center}

\section{分布空间上的社会作用量}

考虑满足连续性方程的 $(m_t,v_t)$：
\[
\partial_tm_t+\nabla\cdot(m_tv_t)=0,
\qquad m(0)=m_0.
\]
势问题的作用量可写成
\focusequation{Wasserstein 空间中的整体优化}{%
\[
\mathcal J(m,v)=
\int_0^T\left[
\int_{\R^d}L(x,v_t(x))m_t(\dd x)+\mathcal F(m_t)
\right]\dd t
+\mathcal G(m_T).
\]}
对约束引入带符号约定的乘子 $-u$，对 $v$ 的变分给出最优反馈，对 $m$ 的变分给出 HJ 方程，约束本身给出连续性方程。于是 MFG 系统恰是这个无限维变分问题的 Euler--Lagrange 条件。

\begin{derivation}[以二次动能为例]
令 $L(v)=\tfrac12\abs v^2$。把 $u$ 作为连续性方程的乘子并对 $v$ 变分：
\[
0=\delta_v\mathcal J
=\int_0^T\int (v+\nabla u)\cdot\delta v\,m\dd x\dd t,
\]
所以 $v^*=-\nabla u$。代回作用量并对 $m$ 变分，得到
\[
-\partial_tu+\frac12\abs{\nabla u}^2
=\frac{\delta\mathcal F}{\delta m}(m,x),
\qquad
\partial_tm-\nabla\cdot(m\nabla u)=0.
\]
\end{derivation}

\begin{pitfall}[势最优解与所有 MFG 均衡并非总是一一对应]
势泛函的全局极小点通常给出均衡，但非凸时还可能有局部极值或鞍点。反过来，一阶 MFG 解只满足驻值条件，不自动保证全局最优。
\end{pitfall}

\section{经验测度：把分布泛函拉回有限维}

给定粒子配置 $\symbf{x}=(x_1,\ldots,x_N)$，定义经验测度与对称函数
\[
m_{\symbf{x}}^N=\frac1N\sum_{i=1}^N\delta_{x_i},
\qquad
F_N(\symbf{x})=\mathcal F(m_{\symbf{x}}^N).
\]
若 $\mathcal F$ 足够光滑，则链式法则给出
\[
D_{x_i}F_N(\symbf{x})=\frac1N
D_m\mathcal F(m_{\symbf{x}}^N,x_i).
\]
二阶导具有不同尺度：对角块通常含 $1/N$ 的“同粒子”项，非对角块是 $1/N^2$ 的“跨粒子”项。示意地，
\[
D_{x_ix_j}^2F_N=
\begin{cases}
O(N^{-1})+O(N^{-2}),&i=j,\\
O(N^{-2}),&i\ne j.
\end{cases}
\]
单个交叉作用很弱，但有 $N(N-1)$ 个交叉块，整体仍可产生有限效应。

\begin{center}
\begin{tikzpicture}[scale=.82]
  \foreach \i in {0,...,6}{
    \foreach \j in {0,...,6}{
      \ifnum\i=\j
        \fill[BlueAccent] (\i*.58,\j*.58) rectangle ++(.48,.48);
      \else
        \fill[PurpleAccent,opacity=.42] (\i*.58,\j*.58) rectangle ++(.48,.48);
      \fi
    }
  }
  \draw[MutedText] (0,0) rectangle (4.0,4.0);
  \node[font=\scriptsize\sffamily,text=BlueAccent,right] at (4.45,3.1) {对角：$O(N^{-1})$};
  \node[font=\scriptsize\sffamily,text=PurpleAccent,right] at (4.45,2.45) {非对角：$O(N^{-2})$};
  \node[font=\scriptsize\sffamily,text=MutedText,right,text width=4.1cm] at (4.45,1.45) {Hessian 的块结构记录“自己移动”与“别人移动”对分布泛函的不同影响。};
\end{tikzpicture}
\captionof{figure}{经验测度泛函 Hessian 的尺度图。}
\end{center}

\section{Hilbert 空间上的 Hamilton--Jacobi 方程}

把分布价值 $\mathcal U(t,m)$ 提升为 $\widetilde{\mathcal U}(t,X)=\mathcal U(t,\Law(X))$。对二次动能，一个典型的无限维 HJ 方程是
\[
\partial_t\widetilde{\mathcal U}(t,X)
-\frac12\norm{D_X\widetilde{\mathcal U}(t,X)}_{L^2}^2
+\widetilde{\mathcal F}(X)=0,
\qquad
\widetilde{\mathcal U}(T,X)=\widetilde{\mathcal G}(X),
\]
符号由时间方向和最小化约定决定。因为提升函数只依赖 $\Law(X)$，它的梯度天然具有 $D_m\mathcal U(m)(X)$ 的形式。

\begin{intuition}[为什么 Hilbert 提升有用]
Wasserstein 空间弯曲且没有线性加法；$L^2$ 空间可以直接使用 Fréchet 导数、半群和经典 HJ 技术。代价是必须始终保持“同分布重排不变”这一对称性。
\end{intuition}

\section{一个相互作用势的几何图像}

令
\[
\mathcal F(m)=\int V(x)m(\dd x)+\frac\kappa2\iint W(x-y)m(\dd x)m(\dd y).
\]
外势 $V$ 塑造地形，相互作用核 $W$ 决定粒子相吸或相斥。$D_m\mathcal F$ 是分布空间中的坡度；若做梯度流，质量沿 $-D_m\mathcal F$ 下坡；若做最优控制/Hamilton 流，还会引入动量和未来价值，因此不能把二者混为一谈。

\begin{center}
\begin{tikzpicture}[x=1cm,y=1cm]
  \draw[axis] (-4.2,0)--(4.2,0) node[right,font=\scriptsize\sffamily] {分布方向};
  \draw[axis] (-3.8,-.2)--(-3.8,3) node[above,font=\scriptsize\sffamily] {$\mathcal F$};
  \draw[BlueAccent,line width=2pt,domain=-3.6:3.8,samples=100] plot (\x,{.10*(\x+2.2)^2*(\x-1.6)^2+.28});
  \fill[YellowAccent] (1.6,.28) circle (3pt);
  \draw[warmflow] (.1,2.0)--(1.25,.55);
  \node[font=\scriptsize\sffamily,text=YellowAccent] at (2.35,.65) {低势能分布};
  \node[font=\scriptsize\sffamily,text=MutedText] at (.4,2.38) {$-\operatorname{grad}_{W_2}\mathcal F$};
\end{tikzpicture}
\captionof{figure}{分布不是沿横轴移动的一个数；图中的“方向”代表整个位移场。能量地形只是无限维几何的二维投影。}
\end{center}

\section{势结构与 Lasry--Lions 单调性}

势性回答“耦合是否来自某个能量”；单调性回答“两个分布之间的耦合是否具有稳定的正向性”。常见条件为
\[
\int_{\R^d}\bigl(F(x,m)-F(x,m')\bigr)(m-m')(\dd x)\ge0.
\]
单调性常用于唯一性，势性常用于存在性和变分构造；二者相关但不等价。

\begin{takeaway}
\begin{enumerate}
  \item 势 MFG 的个体边际成本来自统一泛函 $\mathcal F,\mathcal G$；
  \item MFG 系统是受连续性方程约束的分布作用量的一阶条件；
  \item 经验测度把无限维泛函拉回对称的有限维函数，并产生清晰的 $1/N$ 尺度；
  \item Hilbert 提升把 Wasserstein HJ 方程置于可微的线性空间中。
\end{enumerate}
\end{takeaway}

\begin{exercisebox}
\begin{enumerate}
  \item 对 $\mathcal F(m)=\tfrac12\iint\abs{x-y}^2m(\dd x)m(\dd y)$ 计算 $\delta\mathcal F/\delta m$。
  \item 从 $F_N(\symbf{x})=\mathcal F(m_{\symbf{x}}^N)$ 推导一阶导数的 $1/N$ 因子。
  \item 举例说明“势性但不凸”和“单调但未必显式给出势”的区别。
\end{enumerate}
\end{exercisebox}

\chapter{随机最优控制、动态规划与 HJB 方程}
\lectureribbon{06}{HJB Equations, DPP and Stochastic Optimal Control I}{Andrzej Święch}
\begin{learninggoals}
\begin{itemize}
  \item 建立受控随机微分方程、成本泛函与值函数；
  \item 理解动态规划原理为何是“最优性的时间一致性”；
  \item 用短时间展开与 Itô 公式形式推导 HJB 方程；
  \item 认识 Nisio 半群如何把动态规划写成算子半群。
\end{itemize}
\end{learninggoals}

\section{受控扩散：决策与噪声共同推动状态}

在带滤过的概率空间 $(\Omega,\mathcal F,(\mathcal F_s)_{s\ge0},\Pp)$ 上，令 $W$ 为 $r$ 维 Brown 运动。控制取值于集合 $A$，容许控制 $\alpha_s$ 必须逐步可测，不能偷看未来噪声。状态方程为
\focusequation{受控随机微分方程}{%
\[
\dd X_s=b(s,X_s,\alpha_s)\dd s
+\sigma(s,X_s,\alpha_s)\dd W_s,
\qquad X_t=x.
\]}
漂移 $b$ 是可操纵的平均方向，扩散矩阵 $\sigma$ 决定随机扰动的形状。给定运行成本 $f$ 与终端成本 $g$，
\[
J(t,x;\alpha)=
\E\!\left[
\int_t^T f(s,X_s,\alpha_s)\dd s+g(X_T)
\right],
\qquad
V(t,x)=\inf_{\alpha\in\mathcal A_t}J(t,x;\alpha).
\]
$V(t,x)$ 把从 $(t,x)$ 出发的全部未来随机选择压缩成一个数。

\begin{center}
\begin{tikzpicture}[x=1cm,y=1cm]
  \draw[axis] (-4.4,-1.7)--(4.6,-1.7) node[right,font=\scriptsize\sffamily] {$s$};
  \fill[BlueAccent] (-4,-.2) circle (3.5pt);
  \foreach \k/\c/\yy in {1/BlueAccent/.8,2/PurpleAccent/.2,3/GreenAccent/-.55,4/YellowAccent/1.35,5/OrangeAccent/-1.1}{
    \draw[\c,thick] (-4,-.2)
      .. controls (-2.8,{-.1+.22*\k}) and (-1.6,{\yy-.45}) ..
      (-.6,\yy)
      .. controls (1.0,{\yy+.35-\k*.08}) and (2.6,{\yy-.25+\k*.06}) ..
      (4.2,{\yy+.15});
  }
  \node[font=\scriptsize\sffamily,text=BlueAccent] at (-3.45,.2) {同一初值};
  \node[conceptyellow] at (0,2.35) {控制改变整棵未来轨迹树的分布};
\end{tikzpicture}
\captionof{figure}{噪声产生许多可能路径；策略的目标不是指定一条轨迹，而是优化所有路径上的期望成本。}
\end{center}

\section{动态规划原理：把长问题切成两段}

对 $t\le\tau\le T$，动态规划原理（DPP）写成
\focusequation{Bellman 递归}{%
\[
V(t,x)=\inf_{\alpha\in\mathcal A_t}
\E\!\left[
\int_t^\tau f(s,X_s,\alpha_s)\dd s
+V(\tau,X_\tau)
\right].
\]}
前半段支付即时成本，后半段不再展开所有未来控制，而由 $V(\tau,X_\tau)$ 汇总。它要求最优策略的尾部在到达的随机状态上仍然最优，这就是\keyterm{时间一致性}。

\begin{center}
\begin{tikzpicture}[node distance=12mm and 10mm]
  \node[concept] (now) {现在\\$(t,x)$};
  \node[conceptgreen,right=of now] (tau) {中间状态\\$(\tau,X_\tau)$};
  \node[conceptyellow,right=of tau] (end) {终端\\$g(X_T)$};
  \draw[flow] (now)--node[above,font=\scriptsize\sffamily,text=MutedText]{$\int_t^\tau f\,\dd s$} (tau);
  \draw[warmflow] (tau)--node[above,font=\scriptsize\sffamily,text=MutedText]{$V(\tau,X_\tau)$} (end);
  \draw[softflow,bend right=28] (now) to node[below,font=\scriptsize\sffamily,text=BlueAccent]{一次优化 = 当前一步 + 最优余程} (end);
\end{tikzpicture}
\captionof{figure}{Bellman 原理把连续时间控制变成可无限细分的递归。}
\end{center}

\begin{pitfall}[DPP 不是普通的期望塔性质]
条件期望只能分解一个已给定策略的成本；DPP 还必须证明“前段控制”和“依随机中间状态选择的后段控制”能够可测地拼成容许控制。第八、九讲专门处理这一技术核心。
\end{pitfall}

\section{从 DPP 到 HJB：只看一个极短时间片}

定义固定控制 $a$ 下的扩散生成元
\[
\mathcal L^a\phi(t,x)
=b(t,x,a)\cdot D_x\phi(t,x)
+\frac12\Tr\!\left[\sigma\sigma^\top(t,x,a)D_x^2\phi(t,x)\right].
\]
在 $[t,t+h]$ 暂取常控制 $a$。若 $V$ 光滑，Itô 公式给出
\[
\E[V(t+h,X_{t+h})]-V(t,x)
=\E\int_t^{t+h}\bigl(\partial_sV+\mathcal L^aV\bigr)(s,X_s)\dd s.
\]
代入 DPP，除以 $h$ 并令 $h\downarrow0$，形式上得到
\focusequation{Hamilton--Jacobi--Bellman 方程}{%
\[
\partial_tV(t,x)+
\inf_{a\in A}\left\{
b(t,x,a)\cdot DV
+\frac12\Tr\!\left[\sigma\sigma^\top(t,x,a)D^2V\right]
+f(t,x,a)
\right\}=0,
\quad V(T,x)=g(x).
\]}

\begin{derivation}[每一项来自哪里]
\begin{itemize}
  \item $\partial_tV$：剩余时间缩短造成的价值变化；
  \item $b\cdot DV$：漂移沿价值曲面方向移动；
  \item $\tfrac12\Tr(\sigma\sigma^\top D^2V)$：Brown 运动的二次变差；
  \item $f$：当前瞬时成本；
  \item $\inf_a$：在每个 $(t,x)$ 上重新选择最佳动作。
\end{itemize}
\end{derivation}

\section{Nisio 半群：动态规划的算子形式}

对 $t\le s$ 和终端函数 $\varphi$，定义
\[
(\mathcal S_{t,s}\varphi)(x)
=\inf_{\alpha}
\E\!\left[
\int_t^s f(r,X_r,\alpha_r)\dd r+\varphi(X_s)
\right].
\]
DPP 等价于
\[
\mathcal S_{t,s}\circ\mathcal S_{s,r}
=\mathcal S_{t,r},\qquad t\le s\le r,
\]
并且 $V(t,\cdot)=\mathcal S_{t,T}g$。线性 Markov 半群传播期望，Nisio 半群则在传播中加入最小化，因此通常是非线性的。

\begin{intuition}[HJB 是半群的无穷小生成方程]
把 $s=t+h$，研究 $(\mathcal S_{t,t+h}\varphi-\varphi)/h$ 的极限，得到的正是 HJB 算子。于是 DPP 是有限时间尺度的递归，HJB 是它的微分形式。
\end{intuition}

\section{可算例：一维最小能量控制}

令 $\dd X_s=\alpha_s\dd s$，成本
\[
J(t,x;\alpha)=\int_t^T\frac12\alpha_s^2\dd s+\frac q2X_T^2,
\qquad q>0.
\]
HJB 为
\[
\partial_tV+\inf_a\{aV_x+\tfrac12a^2\}=0
\quad\Longrightarrow\quad
\partial_tV-\frac12(V_x)^2=0.
\]
作二次猜测 $V(t,x)=\tfrac12P(t)x^2$，得到
\[
P'(t)=P(t)^2,\qquad P(T)=q,
\qquad
P(t)=\frac{q}{1+q(T-t)}.
\]
最优反馈为
\[
\alpha^*(t,x)=-V_x(t,x)=-P(t)x.
\]
越靠近终端，$P(t)$ 越大，控制越积极地把状态拉回原点。

\begin{takeaway}
\begin{enumerate}
  \item 值函数把未来所有随机策略的最优期望成本压缩到当前状态；
  \item DPP 表达最优尾策略仍然最优；
  \item 短时间 DPP 加 Itô 公式产生 HJB；
  \item Nisio 半群是 DPP 的算子语言，HJB 是其无穷小生成方程。
\end{enumerate}
\end{takeaway}

\begin{exercisebox}
\begin{enumerate}
  \item 若 $\sigma$ 与控制无关，指出 HJB 中控制最小化只作用于哪些项。
  \item 对 $\dd X_s=\alpha_s\dd s+\eta\dd W_s$ 重做一维二次例子，观察扩散如何改变值函数常数项。
  \item 写出零运行成本时的 Nisio 半群，并与普通 Markov 半群比较。
\end{enumerate}
\end{exercisebox}

\chapter{验证定理与最优反馈控制}
\lectureribbon{07}{HJB Equations, DPP and Stochastic Optimal Control II}{Andrzej Święch}
\begin{learninggoals}
\begin{itemize}
  \item 用 Itô 公式证明任何 HJB 光滑解都是成本的下界；
  \item 理解何时 Hamiltonian 的极小点产生真正的最优反馈；
  \item 区分强控制与弱控制表述，以及退化/一致抛物两种 PDE 情形；
  \item 看清验证定理的适用边界。
\end{itemize}
\end{learninggoals}

\section{HJB 候选解怎样“验证”最优性}

令
\[
\mathcal H(t,x,p,M)
=\inf_{a\in A}\left\{
b(t,x,a)\cdot p
+\frac12\Tr[\sigma\sigma^\top(t,x,a)M]
+f(t,x,a)\right\}.
\]
假设 $v\in C^{1,2}$ 满足
\[
\partial_tv+\mathcal H(t,x,Dv,D^2v)=0,
\qquad v(T,x)=g(x).
\]
对任意容许控制 $\alpha$，由于 infimum 不大于把 $a$ 换成 $\alpha_s$ 后的值，
\[
\partial_tv+\mathcal L^{\alpha_s}v+f(s,x,\alpha_s)\ge0.
\]
沿受控状态使用 Itô 公式并取期望，鞅项消失：
\[
\E[g(X_T)]-v(t,x)
=\E\int_t^T(\partial_sv+\mathcal L^{\alpha_s}v)(s,X_s)\dd s.
\]
于是
\[
J(t,x;\alpha)-v(t,x)
=\E\int_t^T
\bigl(\partial_sv+\mathcal L^{\alpha_s}v+f(s,X_s,\alpha_s)\bigr)\dd s
\ge0.
\]

\begin{theorem}[经典验证定理]
若上述光滑性、增长条件与可积性足以合法使用 Itô 公式，则
\[
v(t,x)\le V(t,x).
\]
若存在容许控制 $\alpha^*$ 几乎处处实现 HJB 中的极小值，则等号成立：
\[
v(t,x)=V(t,x)=J(t,x;\alpha^*).
\]
\end{theorem}

\begin{center}
\begin{tikzpicture}[node distance=12mm and 17mm]
  \node[concept] (hjb) {HJB 光滑解\\$v$};
  \node[conceptyellow,right=of hjb] (ineq) {对任意控制\\$v\le J(\alpha)$};
  \node[conceptgreen,right=of ineq] (attain) {若逐点取到极小\\$J(\alpha^*)=v$};
  \draw[flow] (hjb)--node[above,font=\scriptsize\sffamily,text=MutedText]{Itô} (ineq);
  \draw[warmflow] (ineq)--node[above,font=\scriptsize\sffamily,text=MutedText]{等号条件} (attain);
\end{tikzpicture}
\captionof{figure}{验证证明有两步：先得到普遍下界，再找一条使所有不等式变等式的控制。}
\end{center}

\section{从逐点极小化到反馈控制}

若存在可测选择器
\[
\hat a(t,x)\in\argmin_{a\in A}
\left\{\mathcal L^av(t,x)+f(t,x,a)\right\},
\]
自然候选是 Markov 反馈
\[
\alpha_s^*=\hat a(s,X_s^*),
\]
其中闭环状态满足
\[
\dd X_s^*=b(s,X_s^*,\hat a(s,X_s^*))\dd s
+\sigma(s,X_s^*,\hat a(s,X_s^*))\dd W_s.
\]
但“写出 argmin”还不够：必须确保选择器可测、闭环 SDE 有解、控制确实容许且成本可积。

\begin{center}
\begin{tikzpicture}[node distance=13mm and 18mm]
  \node[concept] (x) {观测状态\\$X_s$};
  \node[conceptyellow,right=of x] (policy) {反馈律\\$\hat a(s,X_s)$};
  \node[conceptgreen,right=of policy] (dyn) {闭环动力学\\$\dd X_s$};
  \draw[flow] (x)--(policy);
  \draw[warmflow] (policy)--(dyn);
  \draw[softflow,bend left=32] (dyn) to node[above,font=\scriptsize\sffamily,text=MutedText]{新状态} (x);
\end{tikzpicture}
\captionof{figure}{反馈不是一条预先画好的轨迹，而是状态—动作—新状态的闭环。}
\end{center}

\section{一个随机线性二次例子}

考虑
\[
\dd X_s=(aX_s+b\alpha_s)\dd s+\eta\dd W_s,
\qquad
J=\E\!\left[\int_t^T\frac12(qX_s^2+r\alpha_s^2)\dd s+\frac h2X_T^2\right],
\]
其中 $r>0$。设
\[
V(t,x)=\frac12P(t)x^2+c(t).
\]
HJB 中关于控制的部分是
\[
b\alpha Px+\frac r2\alpha^2,
\]
故
\[
\alpha^*(t,x)=-\frac{bP(t)}r\,x.
\]
比较 $x^2$ 与常数项得到 Riccati 系统
\[
\begin{aligned}
P'(t)+2aP(t)-\frac{b^2}{r}P(t)^2+q&=0,
&&P(T)=h,\\
c'(t)+\frac{\eta^2}{2}P(t)&=0,
&&c(T)=0.
\end{aligned}
\]
噪声不改变线性反馈斜率的形式，却通过 $c(t)$ 提高期望成本。

\section{一致抛物与退化情形}

若存在 $\lambda>0$ 使
\[
\sigma\sigma^\top(t,x,a)\ge\lambda I
\quad\text{对所有 }(t,x,a),
\]
称方程一致抛物。噪声在所有方向平滑，经典正则性更有希望。若 $\sigma\sigma^\top$ 只有低秩，甚至 $\sigma=0$，HJB 是退化抛物或一阶方程，价值函数更容易出现折角。

\begin{tcolorbox}[colback=NightPanelTwo,colframe=GridLine]
\begin{tabularx}{\textwidth}{@{}>{\sffamily\bfseries\color{BlueAccent}}p{0.22\textwidth}YY@{}}
\toprule
 & \textcolor{GreenAccent}{一致抛物} & \textcolor{YellowAccent}{退化/一阶} \\
\midrule
噪声方向 & 覆盖所有状态方向 & 只覆盖部分方向或没有噪声 \\
平滑作用 & 强 & 弱或缺失 \\
经典解 & 条件合适时可期待 & 常常不能期待 \\
稳健语言 & 经典解与黏性解 & 黏性解通常不可缺少 \\
\bottomrule
\end{tabularx}
\end{tcolorbox}

\section{强表述与弱表述}

\begin{definition}[强控制问题]
概率空间、Brown 运动和滤过预先固定；控制在这套随机基底上选择，状态是给定 Brown 运动的强解。
\end{definition}

\begin{definition}[弱控制问题]
优化时允许概率空间、Brown 运动乃至状态联合分布一并变化，只要求它们满足相应的鞅问题或 SDE 分布约束。
\end{definition}

强表述便于构造路径上的控制；弱表述在扩散系数受控、强解不唯一或需要紧性方法时更灵活。两者是否有相同值，需要额外条件，不能想当然。

\begin{pitfall}[验证定理不负责制造光滑性]
验证定理的逻辑是“若有足够光滑的候选解，则可证明它就是值函数”。它不证明 HJB 必有经典解。现实中值函数常只连续，此时要用黏性解与比较原理替代逐点二阶导数。
\end{pitfall}

\begin{takeaway}
\begin{enumerate}
  \item HJB 光滑解经 Itô 公式给出所有控制成本的下界；
  \item 逐点实现 Hamiltonian 极小值且闭环可解时，反馈控制最优；
  \item 可测选择、闭环适定性和可积性是 argmin 背后的三道门；
  \item 退化扩散会削弱平滑性，推动理论转向黏性解。
\end{enumerate}
\end{takeaway}

\begin{exercisebox}
\begin{enumerate}
  \item 在验证推导中逐行指出哪一步使用了 HJB，哪一步使用了终端条件。
  \item 完成标量 Riccati 例子的系数比较，并说明 $r$ 变大时反馈强度如何变化。
  \item 举出一个逐点 argmin 存在但闭环 SDE 可能不适定的原因。
\end{enumerate}
\end{exercisebox}

\chapter{动态规划背后的可测结构}
\lectureribbon{08}{HJB Equations, DPP and Stochastic Optimal Control III}{Andrzej Święch}
\begin{learninggoals}
\begin{itemize}
  \item 区分适应、逐步可测与可预测控制，理解“不能看未来”的精确定义；
  \item 在典范 Wiener 空间上表示过去路径、未来增量与平移控制；
  \item 理解标准 Borel 空间为何让正则条件概率和可测选择可用；
  \item 为 DPP 证明准备控制拼接与值函数确定性。
\end{itemize}
\end{learninggoals}

\section{为什么最直观的 DPP 还需要一整讲技术}

直观上，到达 $(s,X_s)$ 后“重新选择最优控制”即可。但 $X_s$ 是随机变量，不同样本点需要选择不同的后续控制。若选择方式不可测，拼出的对象就不是合法随机过程，期望也未必有意义。因此必须回答：

\begin{itemize}
  \item 控制究竟能依赖路径的哪一部分？
  \item 给定随机中间状态，如何可测地选近似最优后续控制？
  \item 条件在过去路径上后，未来噪声是否仍是 Brown 运动？
  \item 两段控制拼接后是否仍属于同一容许类？
\end{itemize}

\section{滤过与非预见性}

滤过 $(\mathcal F_s)$ 表示随时间积累的信息。控制 $\alpha_s$ 至少应当适应：$\alpha_s$ 只能使用 $\mathcal F_s$ 中的信息。为了使随机积分与参数化操作稳定，通常要求更强的逐步可测或可预测性质。

\begin{definition}[逐步可测]
过程 $\alpha$ 若对每个 $s$，映射
\[
(r,\omega)\longmapsto\alpha_r(\omega),
\qquad (r,\omega)\in[0,s]\times\Omega
\]
相对于 $\mathcal B([0,s])\otimes\mathcal F_s$ 可测，则称逐步可测。
\end{definition}

\begin{intuition}[信息锥]
在时空图上，时刻 $s$ 的控制只能看到 $s$ 左侧的路径。任何依赖 $W_T-W_s$ 的规则都跨越了信息边界，是预见性控制。
\end{intuition}

\begin{center}
\begin{tikzpicture}[x=1cm,y=1cm]
  \draw[axis] (-4.4,0)--(4.5,0) node[right,font=\scriptsize\sffamily] {时间};
  \draw[very thick,BlueAccent] (-4,0) .. controls (-2.9,.7) and (-1.7,-.45) .. (-.3,.35);
  \draw[very thick,MutedText,dashed] (-.3,.35) .. controls (1.2,1.1) and (2.6,-.8) .. (4.1,.55);
  \draw[YellowAccent,thick] (-.3,-.35)--(-.3,1.45);
  \node[font=\scriptsize\sffamily,text=BlueAccent] at (-2.25,1.15) {可见过去};
  \node[font=\scriptsize\sffamily,text=MutedText] at (2.0,1.15) {不可见未来};
  \node[font=\scriptsize\sffamily,text=YellowAccent] at (-.3,-.7) {$s$};
  \node[conceptpurple] at (0,-1.55) {$\alpha_s$ 可依赖蓝色路径，但不能依赖灰色增量};
\end{tikzpicture}
\captionof{figure}{容许控制的非预见性约束。}
\end{center}

\section{典范 Wiener 空间}

取
\[
\Omega=C([0,T];\R^r),\qquad
W_s(\omega)=\omega(s),
\]
配以 Wiener 测度与自然滤过。每个样本点本身就是一条连续路径，控制可写成路径泛函
\[
\alpha_s(\omega)=\Phi_s\bigl(\omega|_{[0,s]}\bigr).
\]
这一典范表示使“截断过去”“平移未来”和“拼接路径”都变成显式映射，而不是隐藏在抽象概率空间中。

固定确定时刻 $t$，定义未来增量路径
\[
(\theta_t\omega)(r)=\omega(t+r)-\omega(t),\qquad 0\le r\le T-t.
\]
在适当条件下，$\theta_tW$ 与 $\mathcal F_t$ 独立，且仍是从零开始的 Brown 运动。这是把后续问题重启为一个新控制问题的概率基础。

\section{标准 Borel 空间为何重要}

可分完备度量空间及其 Borel $\sigma$-代数称为标准 Borel 空间。欧氏空间、连续路径空间和常见控制路径空间都属于这一类。这个假设看似抽象，却保证两件关键工具：

\begin{enumerate}
  \item 存在良好的正则条件概率，可把“给定过去”表示为概率核；
  \item 多值映射满足条件时存在可测选择，可按随机状态选择近似最优控制。
\end{enumerate}

\begin{pitfall}[逐点选最优不代表可测]
对每个 $x$ 单独说“取一个最优控制 $\alpha^x$”使用了选择公理，却没有保证 $x\mapsto\alpha^x$ 可测。DPP 中需要的不是一堆孤立选择，而是一张可用于随机变量 $X_s$ 的可测控制表。
\end{pitfall}

\section{控制的平移与拼接}

设前段控制为 $\alpha$，后段控制族为 $\beta^y$，其中 $y$ 表示中间状态。确定时刻 $s$ 的拼接控制形式上写成
\[
(\alpha\otimes_s\beta)_r=
\begin{cases}
\alpha_r,&t\le r<s,\\
\beta^{X_s}_r,&s\le r\le T.
\end{cases}
\]
严格构造时，$\beta^{X_s}$ 还要作用于 $s$ 后的增量路径 $\theta_s\omega$。若 $(y,\omega)\mapsto\beta^y(\omega)$ 可测且每个 $\beta^y$ 非预见，则拼接控制仍然容许。

\begin{center}
\begin{tikzpicture}[x=1cm,y=1cm]
  \draw[axis] (-4.3,0)--(4.4,0);
  \draw[BlueAccent,line width=3pt] (-4,0)--(-.35,0);
  \draw[YellowAccent,line width=3pt] (-.35,0)--(4.05,0);
  \fill[SoftWhite] (-.35,0) circle (3pt);
  \node[font=\scriptsize\sffamily,text=BlueAccent] at (-2.1,.55) {前段 $\alpha$};
  \node[font=\scriptsize\sffamily,text=YellowAccent] at (1.8,.55) {按 $X_s$ 选择后段 $\beta^{X_s}$};
  \node[font=\scriptsize\sffamily,text=SoftWhite] at (-.35,-.55) {$s$};
  \draw[softflow] (-.35,1.55)--(-.35,.35);
  \node[conceptpurple] at (-.35,2.05) {状态 $X_s$ 是两段之间唯一必要的 Markov 接口};
\end{tikzpicture}
\captionof{figure}{控制拼接：过去固定，未来根据随机中间状态重新选择。}
\end{center}

\section{为什么值函数是确定函数}

在某些强表述中，先定义的“条件最优成本”可能看起来像 $\mathcal F_t$-可测随机变量。要证明它其实等于常数 $V(t,x)$，通常利用：

\begin{itemize}
  \item 从确定状态 $x$ 出发，未来 Brown 增量与 $\mathcal F_t$ 独立；
  \item 容许控制可在未来增量空间上重表示；
  \item 对过去路径做保持未来 Wiener 测度的变换，值不变。
\end{itemize}

因此 Markov 系数下，过去如何走到 $(t,x)$ 不再影响从该点开始的最优值。若模型含路径依赖，则状态必须扩充为整段历史，值函数也会变成路径泛函。

\section{值函数的连续性估计}

若 $b,\sigma,f,g$ 对 $x$ 一致 Lipschitz，两个同控制、同 Brown 运动驱动的状态满足
\[
\E\sup_{t\le r\le T}\abs{X_r^{t,x,\alpha}-X_r^{t,y,\alpha}}^2
\le C\abs{x-y}^2.
\]
由成本差估计并对控制取 infimum，可得
\[
\abs{V(t,x)-V(t,y)}\le C\abs{x-y}.
\]
时间连续性通常由短时间状态增量 $\E\abs{X_{t+h}-x}\le C\sqrt h$ 推出。连续性既服务于 PDE 理论，也允许在 DPP 证明中用有限状态网格近似随机中间状态。

\begin{takeaway}
\begin{enumerate}
  \item 容许控制必须非预见且具有足够的联合可测性；
  \item 典范 Wiener 空间把过去、未来增量和平移操作显式化；
  \item 标准 Borel 结构保证正则条件概率与可测选择；
  \item 控制拼接是 DPP 的真正技术核心，状态连续性帮助把不可数选择离散化。
\end{enumerate}
\end{takeaway}

\begin{exercisebox}
\begin{enumerate}
  \item 判断 $\alpha_s=\mathbf 1_{\{W_T>0\}}$ 是否容许，并说明原因。
  \item 在典范路径空间写出 $\theta_t$，验证它只记录 $t$ 后的增量。
  \item 说明为什么“每个 $x$ 有一个 $\varepsilon$-最优控制”还不足以直接代入随机变量 $X_s$。
\end{enumerate}
\end{exercisebox}

\chapter{动态规划原理的严格证明}
\lectureribbon{09}{HJB Equations, DPP and Stochastic Optimal Control IV}{Andrzej Święch}
\begin{learninggoals}
\begin{itemize}
  \item 把 DPP 拆成两个方向相反的不等式，并辨认各自的难点；
  \item 使用条件概率、Brown 运动增量与 SDE 流性质处理后续成本；
  \item 用状态分区和 $\varepsilon$-最优控制完成可测拼接；
  \item 理解哪些正则性假设负责控制误差、哪些假设负责可测性。
\end{itemize}
\end{learninggoals}

\section{要证明的等式}

对确定中间时刻 $s\in[t,T]$，记
\[
\mathcal G_{t,s}^{\alpha}[\varphi](x)
=\E\!\left[
\int_t^s f(r,X_r^{t,x,\alpha},\alpha_r)\dd r
+\varphi(X_s^{t,x,\alpha})
\right].
\]
动态规划原理可紧凑地写为
\focusequation{DPP 的算子形式}{%
\[
V(t,x)=\inf_{\alpha\in\mathcal A_t}
\mathcal G_{t,s}^{\alpha}[V(s,\cdot)](x).
\]}
等号左边允许在整个区间上选择一条控制；右边只显式选择前段控制，把后段最优值压缩为 $V(s,X_s)$。证明要说明这两种优化方式能力完全相同。

\begin{center}
\begin{tikzpicture}[node distance=11mm and 16mm]
  \node[concept] (whole) {一次优化整段\\$V(t,x)$};
  \node[conceptyellow,right=of whole] (split) {前段控制\\$[t,s]$};
  \node[conceptgreen,right=of split] (tail) {随机状态后的最优余程\\$V(s,X_s)$};
  \draw[flow] (whole)--node[above,font=\scriptsize\sffamily,text=MutedText]{分解} (split);
  \draw[warmflow] (split)--node[above,font=\scriptsize\sffamily,text=MutedText]{拼接} (tail);
  \draw[softflow,bend right=30] (tail) to node[below,font=\scriptsize\sffamily,text=BlueAccent]{恢复整段策略} (whole);
\end{tikzpicture}
\captionof{figure}{DPP 证明的两次旅行：先把整段策略切开，再把依状态选择的后段策略拼回去。}
\end{center}

\section{状态流性质：在中间时刻重新启动}

在 Lipschitz 条件下，受控 SDE 有路径唯一解。若控制在 $s$ 后平移为 $\alpha^{s,\omega}$，则几乎处处
\[
X_r^{t,x,\alpha}(\omega)
=X_r^{s,X_s^{t,x,\alpha}(\omega),\alpha^{s,\omega}}
\bigl(\theta_s\omega\bigr),
\qquad r\ge s.
\]
这条式子不是说两边使用完全相同的原始路径，而是说：给定过去后，后续状态由中间状态、平移后的控制和新的 Brown 增量唯一决定。这是 Markov 重启的路径版。

\section{第一方向：整段成本不小于“前段 + 最优余程”}

任取整个区间上的容许控制 $\alpha$。把成本在 $s$ 处分开：
\[
J(t,x;\alpha)
=\E\!\left[
\int_t^s f(r,X_r,\alpha_r)\dd r
+\E\!\left[
\left.
\int_s^T f(r,X_r,\alpha_r)\dd r+g(X_T)
\right|\mathcal F_s
\right]\right].
\]
在给定 $\mathcal F_s$ 后，$\alpha$ 的尾部只是从随机状态 $X_s$ 出发的一种容许后续控制，因此条件后续成本至少为最优值：
\[
\E\!\left[
\left.
\int_s^T f\,\dd r+g(X_T)
\right|\mathcal F_s
\right]\ge V(s,X_s).
\]
于是
\[
J(t,x;\alpha)
\ge\mathcal G_{t,s}^{\alpha}[V(s,\cdot)](x).
\]
再对整个区间控制取 infimum，得到
\focusequation{DPP 的“$\ge$”方向}{%
\[
V(t,x)\ge
\inf_{\alpha}\mathcal G_{t,s}^{\alpha}[V(s,\cdot)](x).
\]}

\begin{intuition}[这一方向为什么相对容易]
任何既定尾策略都不会优于尾部问题的最优值。只要条件化后尾控制仍然合法，结论便直接来自“最优值 $\le$ 任意可行值”。
\end{intuition}

\section{第二方向：把近似最优余程可测地接上}

反向不等式不能简单地写“到 $X_s$ 后取最优控制”，因为 $X_s$ 取不可数多个值，而且最优控制未必存在。标准办法是“离散化状态 + 选择近似最优控制 + 拼接”。

\subsection{第一步：为代表点选控制}

把状态空间分成直径不超过 $\delta$ 的 Borel 集合 $B_k$，在每个非空单元选代表点 $y_k$。按值函数定义，存在控制 $\beta^k$ 使
\[
J(s,y_k;\beta^k)\le V(s,y_k)+\varepsilon.
\]

\subsection{第二步：按随机状态选择}

定义后段控制
\[
\beta_r(\omega)
=\sum_k\mathbf 1_{\{X_s(\omega)\in B_k\}}\,
\beta_r^k(\theta_s\omega),
\qquad r\ge s.
\]
事件 $\{X_s\in B_k\}$ 属于 $\mathcal F_s$，而 $\beta^k$ 只使用 $s$ 后增量，所以 $\beta$ 仍非预见且可测。

\subsection{第三步：用稳定性控制网格误差}

若值函数和成本对初值 Lipschitz，则对 $X_s\in B_k$，
\[
J(s,X_s;\beta^k)
\le J(s,y_k;\beta^k)+C\abs{X_s-y_k}
\le V(s,X_s)+\varepsilon+C'\delta.
\]
把任意前段控制 $\alpha$ 与这条后段控制拼接，得到
\[
V(t,x)
\le\mathcal G_{t,s}^{\alpha}[V(s,\cdot)](x)
+\varepsilon+C'\delta.
\]
先对前段 $\alpha$ 取 infimum，再令 $\varepsilon,\delta\downarrow0$：
\focusequation{DPP 的“$\le$”方向}{%
\[
V(t,x)\le
\inf_{\alpha}\mathcal G_{t,s}^{\alpha}[V(s,\cdot)](x).
\]}

\begin{center}
\begin{tikzpicture}[scale=.92]
  \draw[guide] (-4.3,-1.8) grid[xstep=1.1,ystep=.9] (4.5,1.8);
  \fill[YellowAccent] (-1.0,.55) circle (4pt);
  \node[font=\scriptsize\sffamily,text=YellowAccent] at (-.35,.92) {$X_s$};
  \foreach \x/\y/\lab in {-3.75/1.35/1,-2.65/.45/2,-1.55/1.35/3,-.45/.45/4,.65/1.35/5,1.75/.45/6,2.85/1.35/7,3.95/.45/8}{
    \fill[BlueAccent] (\x,\y) circle (2.5pt);
  }
  \draw[warmflow] (-1,.55)--(-1.55,1.35);
  \node[conceptpurple] at (0,-2.55) {找到 $X_s$ 所在网格，调用该格代表点的 $\varepsilon$-最优尾控制};
\end{tikzpicture}
\captionof{figure}{状态分区把不可数的控制选择转成可数个可测分支；连续性把空间离散误差压到零。}
\end{center}

\section{两方向使用的工具并不相同}

\begin{tcolorbox}[colback=NightPanelTwo,colframe=GridLine]
\begin{tabularx}{\textwidth}{@{}>{\sffamily\bfseries\color{BlueAccent}}p{0.2\textwidth}YY@{}}
\toprule
 & \textcolor{GreenAccent}{$V\ge$ 方向} & \textcolor{YellowAccent}{$V\le$ 方向} \\
\midrule
核心动作 & 条件化既定控制 & 构造一条新拼接控制 \\
最优性使用 & 尾成本 $\ge$ 尾部最优值 & 每格选 $\varepsilon$-最优尾控制 \\
主要工具 & 正则条件概率、状态流性质 & Borel 分区、可测选择、稳定性 \\
主要风险 & 条件后的尾控制不在容许类 & 点态选择不可测或拼接失效 \\
\bottomrule
\end{tabularx}
\end{tcolorbox}

\section{从 DPP 再到 Nisio 半群}

定义非线性算子
\[
(\mathcal S_{t,s}\varphi)(x)
=\inf_{\alpha}\mathcal G_{t,s}^{\alpha}[\varphi](x).
\]
刚证明的 DPP 等价于
\[
V(t,\cdot)=\mathcal S_{t,s}V(s,\cdot),
\]
而对任意终端函数 $\varphi$，同样的切分与拼接论证给出
\[
\mathcal S_{t,s}\mathcal S_{s,r}\varphi
=\mathcal S_{t,r}\varphi.
\]
因此半群性质不是形式记号，而是完整 DPP 证明在算子层面的压缩。

\begin{pitfall}[最优控制可能不存在，但 DPP 仍可成立]
值函数由 infimum 定义，未必有控制真正取到它。证明中始终使用 $\varepsilon$-最优控制，最后令 $\varepsilon\downarrow0$；这避免了不必要的紧性或选择假设。
\end{pitfall}

\begin{takeaway}
\begin{enumerate}
  \item “$\ge$”方向：条件化任意整段控制，尾成本不小于尾部最优值；
  \item “$\le$”方向：按中间状态可测地拼接近似最优尾控制；
  \item 状态连续性负责网格误差，可测结构负责拼接合法性；
  \item 两个方向合并后，得到 Bellman/Nisio 半群。
\end{enumerate}
\end{takeaway}

\begin{exercisebox}
\begin{enumerate}
  \item 不看正文，独立写出 DPP 两个不等式各自从哪一个控制出发。
  \item 若 $V(s,\cdot)$ 只有一致连续性而非 Lipschitz，如何修改网格误差估计？
  \item 解释为何 $\beta^{X_s}$ 必须使用 $s$ 后 Brown 增量，而不能重新使用过去噪声。
\end{enumerate}
\end{exercisebox}

\chapter{从动态规划到黏性解}
\lectureribbon{10}{HJB Equations, DPP and Stochastic Optimal Control V}{Andrzej Święch}
\begin{learninggoals}
\begin{itemize}
  \item 把确定时刻 DPP 推广到停止时刻；
  \item 理解测试函数“从上接触/从下接触”如何替代不存在的导数；
  \item 从 DPP 推导 HJB 的黏性次解与上解不等式；
  \item 用比较原理完成“值函数是唯一黏性解”的闭环。
\end{itemize}
\end{learninggoals}

\section{停止时刻版动态规划}

若 $\tau$ 是取值于 $[t,T]$ 的停止时刻，则强 Markov 思想给出
\focusequation{停止时刻 DPP}{%
\[
V(t,x)=\inf_{\alpha}
\E\!\left[
\int_t^\tau f(s,X_s,\alpha_s)\dd s
+V(\tau,X_\tau)
\right].
\]}
证明通常分两步：先对只取有限多个值的简单停止时刻逐层拼接，再用 $\tau_n\downarrow\tau$、状态估计和值函数连续性取极限。停止时刻 DPP 允许在“首次离开小邻域”时停止，是黏性解局部证明的关键。

\begin{examplebox}[退出时刻]
取圆柱邻域
\[
Q=(t_0,t_0+h)\times B_\rho(x_0),
\]
令 $\tau$ 为过程首次离开 $Q$ 的时刻与 $t_0+h$ 的较小者。测试函数只需在 $Q$ 内控制值函数，停止时刻保证 Itô 展开不会跑出这个局部区域。
\end{examplebox}

\section{为什么经典导数不够}

即使 $b,\sigma,f,g$ 光滑，infimum 也可能使值函数出现折角：不同控制在状态空间的切换面上给出相同价值，但两侧梯度不同。HJB 需要二阶导数，值函数却可能只有连续性。黏性解的做法是：

\begin{quote}
\centering
\emterm{不对粗糙的 $V$ 求导；只对与它相切的光滑函数 $\phi$ 求导。}
\end{quote}

定义
\[
\mathscr H(t,x,p,M)
=\sup_{a\in A}\left\{
-b(t,x,a)\cdot p
-\frac12\Tr[\sigma\sigma^\top(t,x,a)M]
-f(t,x,a)\right\}.
\]
HJB 写成适合比较理论的形式
\focusequation{HJB 的 proper 形式}{%
\[
-\partial_tV+\mathscr H(t,x,DV,D^2V)=0,
\qquad V(T,x)=g(x).
\]}
它与第六讲的 $\partial_tV+\inf_a\{\mathcal L^aV+f\}=0$ 完全等价。

\section{从上接触与从下接触}

\begin{definition}[黏性次解]
上半连续函数 $u$ 是次解，如果每当 $\phi\in C^{1,2}$ 使 $u-\phi$ 在 $(t_0,x_0)$ 取得局部极大值，就有
\[
-\partial_t\phi(t_0,x_0)
+\mathscr H(t_0,x_0,D\phi,D^2\phi)\le0.
\]
\end{definition}

\begin{definition}[黏性上解]
下半连续函数 $w$ 是上解，如果每当 $\phi\in C^{1,2}$ 使 $w-\phi$ 在 $(t_0,x_0)$ 取得局部极小值，就有
\[
-\partial_t\phi(t_0,x_0)
+\mathscr H(t_0,x_0,D\phi,D^2\phi)\ge0.
\]
\end{definition}

连续函数同时为次解和上解，便是黏性解。

\begin{center}
\begin{tikzpicture}[x=1cm,y=1cm]
  \draw[axis] (-4.2,0)--(4.3,0) node[right,font=\scriptsize\sffamily] {$x$};
  \draw[axis] (0,-.4)--(0,3.2);
  \draw[BlueAccent,line width=2pt] plot[smooth] coordinates {(-3.7,.5)(-2.7,1.25)(-1.7,.7)(-.7,1.45)(0,1.2)(.7,1.55)(1.5,.85)(2.5,1.35)(3.7,.65)};
  \draw[YellowAccent,very thick,domain=-2.3:2.3,samples=60] plot (\x,{1.2+.12*(\x)^2});
  \draw[GreenAccent,very thick,domain=-2.1:2.1,samples=60] plot (\x,{1.2-.11*(\x)^2});
  \fill[SoftWhite] (0,1.2) circle (2.8pt);
  \node[font=\scriptsize\sffamily,text=YellowAccent] at (2.8,2.0) {上方测试函数};
  \node[font=\scriptsize\sffamily,text=GreenAccent] at (2.7,.42) {下方测试函数};
  \node[font=\scriptsize\sffamily,text=BlueAccent] at (-3.1,1.75) {$V$ 可有折角};
\end{tikzpicture}
\captionof{figure}{测试函数在接触点借出自己的导数；它无需在整个区域逼近值函数。}
\end{center}

\section{DPP 推出次解不等式}

设 $\phi$ 从上方接触 $V$ 于 $(t_0,x_0)$，即局部有 $V\le\phi$ 且接触点相等。固定任意动作 $a$，在一个短时间/退出时刻前使用常控制 $a$。由 DPP 的“值不大于任意可行策略成本”，
\[
V(t_0,x_0)
\le\E\!\left[\int_{t_0}^{\tau}f(s,X_s,a)\dd s+V(\tau,X_\tau)\right]
\le\E\!\left[\int_{t_0}^{\tau}f(s,X_s,a)\dd s+\phi(\tau,X_\tau)\right].
\]
对 $\phi$ 使用 Itô 公式、除以预期时间并让邻域缩小，得到对每个 $a$ 都有
\[
\partial_t\phi+\mathcal L^a\phi+f(\cdot,a)\ge0.
\]
等价地，
\[
-\partial_t\phi
+\sup_a\{-\mathcal L^a\phi-f(\cdot,a)\}\le0,
\]
这正是次解不等式。

\section{DPP 推出上解不等式}

若 $\phi$ 从下方接触 $V$，则局部 $V\ge\phi$。取短区间上的 $\varepsilon$-最优控制 $\alpha^\varepsilon$，由 DPP
\[
V(t_0,x_0)+\varepsilon h
\ge\E\!\left[\int_{t_0}^{\tau}f(s,X_s,\alpha_s^\varepsilon)\dd s
+V(\tau,X_\tau)\right]
\ge\E\!\left[\int_{t_0}^{\tau}f\,\dd s+\phi(\tau,X_\tau)\right].
\]
Itô 展开并令 $\varepsilon,h\downarrow0$，得到
\[
-\partial_t\phi+\mathscr H(t_0,x_0,D\phi,D^2\phi)\ge0.
\]
次解使用\blue{任意固定控制}，上解使用\yellow{近似最优控制}；这与第九讲 DPP 两个方向的结构完全呼应。

\begin{center}
\begin{tikzpicture}[node distance=12mm and 17mm]
  \node[concept] (dpp) {动态规划原理};
  \node[conceptyellow,right=of dpp] (sub) {上方接触\\次解不等式};
  \node[conceptgreen,below=of sub] (super) {下方接触\\上解不等式};
  \node[conceptpurple,right=15mm of $(sub)!0.5!(super)$] (unique) {比较原理\\唯一连续解};
  \draw[flow] (dpp)--node[above,font=\scriptsize\sffamily,text=MutedText]{任意控制} (sub);
  \draw[warmflow] (dpp)--node[below,font=\scriptsize\sffamily,text=MutedText]{$\varepsilon$-最优控制} (super);
  \draw[softflow] (sub)--(unique);
  \draw[softflow] (super)--(unique);
\end{tikzpicture}
\captionof{figure}{概率控制与 PDE 理论的闭环。}
\end{center}

\section{比较原理：唯一性的发动机}

典型比较结论是：若 $u$ 是适当增长的上半连续次解，$w$ 是下半连续上解，且
\[
u(T,x)\le w(T,x),
\]
则在整个区域 $u\le w$。证明常使用 doubling of variables：
\[
\sup_{t,x,y}\left\{
u(t,x)-w(t,y)-\frac{\abs{x-y}^2}{2\varepsilon}
\right\},
\]
迫使最大点满足 $x\approx y$，再用 Crandall--Ishii 型矩阵不等式比较两边测试二阶导。比较原理带来三件事：

\begin{enumerate}
  \item 黏性解唯一；
  \item 由 DPP 得到的值函数就是该唯一解；
  \item 近似模型和数值格式只要保持单调稳定，就有希望收敛到同一对象。
\end{enumerate}

\section{不可微但完全正确的值函数}

考虑确定性系统 $\dot X=a$、$\abs a\le1$，无运行成本，终端成本 $g(x)=\abs x$。剩余时间内最多向原点移动 $T-t$，故
\[
V(t,x)=\max\{\abs x-(T-t),0\}.
\]
它满足
\[
\partial_tV-\abs{V_x}=0,\qquad V(T,x)=\abs x,
\]
却在 $\abs x=T-t$ 处不可微。经典解语言会失败；测试函数仍能在折角处给出正确的不等式，所以 $V$ 是黏性解。

\begin{center}
\begin{tikzpicture}[x=1cm,y=1cm]
  \draw[axis] (-4.2,0)--(4.2,0) node[right,font=\scriptsize\sffamily] {$x$};
  \draw[axis] (0,-.2)--(0,2.8) node[above,font=\scriptsize\sffamily] {$V$};
  \draw[BlueAccent,line width=2.2pt] (-3.8,2.15)--(-1.35,0)--(1.35,0)--(3.8,2.15);
  \fill[YellowAccent] (-1.35,0) circle (3pt);
  \fill[YellowAccent] (1.35,0) circle (3pt);
  \node[font=\scriptsize\sffamily,text=YellowAccent] at (0,.55) {可达原点：成本为 0};
  \node[font=\scriptsize\sffamily,text=MutedText] at (2.7,1.25) {来不及完全回到原点};
\end{tikzpicture}
\captionof{figure}{控制切换制造折角；折角是问题结构，不是解的缺陷。}
\end{center}

\section{稳定性与无限维延伸}

黏性解对局部一致收敛稳定：若 HJB 算子和解序列合适收敛，极限仍是极限方程的黏性解。这对离散化、噪声消失极限和有限玩家极限至关重要。

在 Hilbert 空间中，状态 $X$ 可代表随机变量、场或延迟系统；HJB 形式仍是
\[
-\partial_tV+\mathscr H(t,X,DV,D^2V)=0.
\]
困难在于二阶算子的迹可能无穷、紧性变弱、测试函数类别需要调整。尽管技术升级，核心仍是 DPP、接触测试与比较。

\begin{takeaway}
\begin{enumerate}
  \item 停止时刻 DPP 把局部退出问题纳入 Bellman 递归；
  \item 测试函数借出导数，使连续但不可微的值函数仍满足 HJB；
  \item 次解来自任意控制，上解来自 $\varepsilon$-最优控制；
  \item 比较原理给出唯一性，并把控制值函数与 PDE 解完全识别。
\end{enumerate}
\end{takeaway}

\begin{exercisebox}
\begin{enumerate}
  \item 用简单停止时刻 $\tau=\sum_k t_k\mathbf 1_{A_k}$ 写出 DPP 的分层拼接。
  \item 对 $V(t,x)=\max\{\abs x-(T-t),0\}$ 分区域验证 HJB，并说明折角处为何要用黏性定义。
  \item 在次解证明中解释为什么从上接触给出 $V(\tau,X_\tau)\le\phi(\tau,X_\tau)$。
\end{enumerate}
\end{exercisebox}


\appendix
\chapter{符号表与延伸阅读}
\section{十讲速查}

\begin{tcolorbox}[colback=NightPanelTwo,colframe=GridLine]
\small
\begin{tabularx}{\textwidth}{@{}>{\sffamily\bfseries\color{BlueAccent}}p{0.08\textwidth}>{\RaggedRight\arraybackslash}p{0.29\textwidth}Y@{}}
\toprule
讲次 & 主题 & 一句话核心 \\
\midrule
01 & 从 $N$ 人博弈到 MFG & 个体最优与总体分布构成前向--后向不动点。\\
02 & Wasserstein 几何 & 分布的距离、速度和梯度都可由最优搬运定义。\\
03 & Lions 导数与主方程 & 在 Hilbert 空间求导，再降回 law-invariant 的分布变量。\\
04 & 特征线与流 & 粒子轨迹的推前等价于连续性方程，流满足时间拼接。\\
05 & 势平均场博弈 & 个体边际成本来自同一分布势能的变分。\\
06 & 随机控制与 HJB & DPP 的短时间极限加 Itô 公式产生 HJB。\\
07 & 验证与反馈 & HJB 解给普遍下界，逐点极小反馈使下界取等。\\
08 & 可测结构 & 条件概率、标准 Borel 与控制拼接让直觉 DPP 合法。\\
09 & DPP 证明 & 条件化给一个方向，可测的 $\varepsilon$-最优拼接给另一个方向。\\
10 & 黏性解 & 测试函数替代缺失导数，比较原理给出唯一性。\\
\bottomrule
\end{tabularx}
\end{tcolorbox}

\section{常用符号}

\begin{tcolorbox}[colback=NightPanelTwo,colframe=GridLine]
\begin{tabularx}{\textwidth}{@{}>{\color{YellowAccent}}p{0.23\textwidth}Y@{}}
\toprule
\textcolor{BlueAccent}{符号} & \textcolor{BlueAccent}{含义} \\
\midrule
$\PP_2(\R^d)$ & 具有有限二阶矩的 Borel 概率测度空间 \\
$\Gamma(\mu,\nu)$ & 边缘分布为 $\mu,\nu$ 的所有耦合 \\
$W_2(\mu,\nu)$ & 二阶 Wasserstein 距离 \\
$T_\#\mu$ & 映射 $T$ 对测度 $\mu$ 的推前 \\
$\Law(X)$ & 随机变量 $X$ 的分布 \\
$D_m\mathcal F(m,x)$ & 分布泛函的 Lions/Wasserstein 向量导数 \\
$\delta\mathcal F/\delta m$ & 标量线性泛函导数，只确定到常数 \\
$m_t^N$ & $N$ 个粒子的经验测度 \\
$u(t,x)$ & 给定均衡分布流上的个体价值 \\
$U(t,x,m)$ & 同时依赖个体状态和总体分布的主方程解 \\
$V(t,x)$ & 随机最优控制值函数 \\
$\mathcal L^a$ & 固定动作 $a$ 的扩散生成元 \\
$\mathcal S_{t,s}$ & Bellman--Nisio 非线性半群 \\
\bottomrule
\end{tabularx}
\end{tcolorbox}

\section{符号约定对照}

本讲义的随机控制采用最小化成本：
\[
V(t,x)=\inf_\alpha\E\!\left[\int_t^Tf\,\dd s+g(X_T)\right].
\]
因此 HJB 有两种完全等价的写法：
\[
\partial_tV+\inf_a\{\mathcal L^aV+f^a\}=0
\quad\Longleftrightarrow\quad
-\partial_tV+\sup_a\{-\mathcal L^aV-f^a\}=0.
\]
经典计算常用左式；黏性比较理论常用右式。若资料采用最大化收益、反向时间或 $H(x,p)=\sup_a\{a\cdot p-L\}$，若干符号会翻转，但 Bellman 结构不变。

\begin{pitfall}[核对公式的四步]
\begin{enumerate}
  \item 问题是最小化成本还是最大化收益？
  \item 终端条件在 $T$，还是初始条件在 $0$？
  \item SDE 扩散写成 $\sigma\dd W$ 还是 $\sqrt{2\nu}\dd W$？
  \item Hamiltonian 对速度/控制的 Legendre 变换用了哪个正负号？
\end{enumerate}
\end{pitfall}

\section{核心公式卡}

\begin{multicols}{2}
\begin{definition}[最优搬运]
\[
W_2^2(\mu,\nu)=
\inf_{\gamma\in\Gamma(\mu,\nu)}
\int\abs{x-y}^2\dd\gamma.
\]
\end{definition}

\begin{definition}[连续性方程]
\[
\partial_tm+\nabla\cdot(mv)=0.
\]
\end{definition}

\begin{definition}[分布链式法则]
\[
\frac{\dd}{\dd t}\mathcal F(m_t)
=\int D_m\mathcal F(m_t,x)\cdot v_t(x)m_t(\dd x).
\]
\end{definition}

\begin{definition}[Hilbert 提升]
\[
\begin{aligned}
\widetilde{\mathcal F}(X)&=\mathcal F(\Law(X)),\\
D\widetilde{\mathcal F}(X)&=D_m\mathcal F(\Law X,X).
\end{aligned}
\]
\end{definition}

\begin{definition}[受控扩散]
\[
\dd X_s=b(s,X_s,\alpha_s)\dd s
+\sigma(s,X_s,\alpha_s)\dd W_s.
\]
\end{definition}

\begin{definition}[动态规划]
\[
V(t,x)=\inf_\alpha\E\!\left[
\int_t^\tau f\,\dd s+V(\tau,X_\tau)\right].
\]
\end{definition}
\end{multicols}

\section{一张总关系图}

\begin{center}
\begin{tikzpicture}[node distance=10mm and 13mm,scale=.94,transform shape]
  \node[concept] (nash) {$N$ 人 Nash 系统};
  \node[conceptpurple,right=of nash] (emp) {经验测度\\$m^N$};
  \node[conceptgreen,right=of emp] (mfg) {MFG 系统\\$u,m$};
  \node[conceptyellow,right=of mfg] (master) {主方程\\$U(t,x,m)$};
  \draw[flow] (nash)--node[above,font=\scriptsize\sffamily,text=MutedText]{$N\to\infty$} (emp);
  \draw[flow] (emp)--(mfg);
  \draw[warmflow] (mfg)--node[above,font=\scriptsize\sffamily,text=MutedText]{所有初值统一化} (master);

  \node[concept,below=15mm of nash] (sde) {受控 SDE};
  \node[conceptpurple,right=of sde] (value) {值函数 $V$};
  \node[conceptgreen,right=of value] (dpp) {动态规划 DPP};
  \node[conceptyellow,right=of dpp] (hjb) {HJB 黏性解};
  \draw[flow] (sde)--(value);
  \draw[flow] (value)--(dpp);
  \draw[warmflow] (dpp)--(hjb);

  \draw[softflow,dashed] (emp)--node[right,font=\scriptsize\sffamily,text=MutedText]{Wasserstein 几何} (value);
  \draw[softflow,dashed] (master)--node[right,font=\scriptsize\sffamily,text=MutedText]{无限维 Bellman} (hjb);
\end{tikzpicture}
\captionof{figure}{上半部是群体博弈极限，下半部是单体随机控制；两者由分布几何和 Bellman 原理连接。}
\end{center}

\section{课程来源与重构原则}

本讲义对应 Tohoku Forum for Creativity 2017 课程《Optimal Control and Partial Differential Equations》十场讲座。第一至第五讲由 Wilfrid Gangbo 主讲，主题为平均场博弈方程；第六至第十讲由 Andrzej Święch 主讲，主题为 HJB 方程、动态规划原理与随机最优控制。

正文根据完整中文精校字幕重构：保留课程论证主线，统一中文术语，补足跨讲依赖，并用向量示意图解释空间关系。它是学习型讲义，不代替原讲座中的全部技术假设与证明细节。

\section{建议的延伸阅读路径}

\begin{enumerate}
  \item \emterm{最优传输与 Wasserstein 几何}：先系统学习最优耦合、测地线、连续性方程与梯度流；
  \item \emterm{平均场博弈}：再学习 MFG 系统的存在唯一性、Lasry--Lions 单调性与主方程；
  \item \emterm{随机控制}：掌握 SDE、Itô 公式、DPP、验证定理和强/弱表述；
  \item \emterm{黏性解}：重点学习半连续包、测试函数、比较原理与稳定性；
  \item \emterm{无限维控制}：最后把状态推广到 Hilbert 空间与概率测度空间。
\end{enumerate}

\begin{takeaway}[全课程最后一页]
\[
\boxed{
\text{局部最优选择}
\ \xrightarrow{\ \text{动态规划}\ }\ 
\text{价值方程}
\ \xrightarrow{\ \text{反馈}\ }\ 
\text{状态/分布演化}
}
\]
\[
\boxed{
\text{有限玩家}
\ \xrightarrow{\ N\to\infty\ }\ 
\text{平均场博弈}
\ \xrightarrow{\ \text{分布微分}\ }\ 
\text{主方程}
}
\]
两条链的共同语言是：\blue{状态空间的几何}、\yellow{未来价值的递归}、\green{最优反馈的闭环}。
\end{takeaway}


\backmatter
\chapter*{结语：同一条 Bellman 主线}
\addcontentsline{toc}{chapter}{结语：同一条 Bellman 主线}

从单个受控粒子到无穷多个相互作用的理性主体，课程始终在回答同一问题：\emterm{今天的最佳选择，如何由明天的最优价值决定？} 有限维中，这个递归关系是动态规划与 HJB 方程；分布空间中，它演化为平均场博弈系统与主方程。对象变了，Bellman 思想没有变。

\begin{center}
\begin{tikzpicture}[node distance=12mm]
\node[concept] (local) {局部决策\\控制 $a$};
\node[conceptyellow,right=of local] (value) {全局价值\\$V$ 或 $U$};
\node[conceptgreen,right=of value] (pde) {偏微分方程\\HJB / 主方程};
\draw[flow] (local)--(value);
\draw[warmflow] (value)--(pde);
\draw[softflow,bend left=35] (pde) to node[above,font=\scriptsize\sffamily,text=MutedText]{反馈最优策略} (local);
\end{tikzpicture}
\end{center}

\vfill
\begin{center}
\sffamily\small\textcolor{MutedText}{由中文精校字幕重构 · 2026 年 8 月}
\end{center}

\end{document}
